\section{Methods: experimental setup}
% [notes moved to thesis_notes_moved.md]
% CODE INPUT
This section sets out how every experiment in this chapter is run. ``Datasets and splitting'' describes the 16 datasets from six domains and how their folds are cut. ``Representations by domain'' lists the base representations offered to the search in each domain. ``Models and metrics'' gives the downstream models and the metric used on each dataset. ``Comparators, controls and ablations'' describes the experiments that test whether the hybrid's gain is real. The last subsection sets up three further studies: reusing search settings on a second dataset, antibody mutational scans, and data held at separate sites.

% CODE INPUT
Every experiment reuses one protocol. Each of the 16 datasets is split into five outer folds, and no group is split across folds. One analysis of the antibody expression scan uses three outer folds of unequal size instead; this is stated where that result is reported. The hybrid search sees only the training samples of each outer fold. The selected features are then rebuilt on those training samples and applied unchanged to the held-out samples.

% CODE INPUT
The hybrid is compared against the best single base representation, used at its full number of features. This asks whether the hybrid beats what a user would otherwise run. The comparator is chosen from the search's own single-representation scores on the training samples of the outer fold. It is then refit on those training samples and scored on the same held-out samples as the hybrid. It is chosen separately for each downstream model and each fold, so it can be a different representation in each fold. Choosing it on the held-out samples would give it an advantage the hybrid does not get. The hybrid is held to its budget of base features, with any augmentation features added on top. The comparator is held to no budget, so it is often many times wider than the hybrid. Both widths are recorded for every fold. A difference between the two is therefore partly a difference in width, not only in method. The main quantity per fold is the hybrid's score minus the best single representation's score. It is signed so that a positive value always means the hybrid scored better; for RMSE the sign is flipped. For each dataset and downstream model, we report the per-fold scores of the hybrid and of every single representation, and their spread across the five folds.

% CODE INPUT
Chapter~\ref{chap:similarity} reuses this protocol with noise added to the training labels of each outer fold. The held-out labels stay clean.


\subsection{Datasets and splitting}
% [notes moved to thesis_notes_moved.md]

% TODO: review, needs manual review
The hybrid experiments use 16 datasets from six domains: five small-molecule ADME and toxicity endpoints, one quantum-chemical property, two reaction-yield sets, four antibody sets, two nucleotide-sequence tasks and two cancer classification tasks from TCGA (Table~\ref{tab:datasets}). Every score in this chapter is computed on held-out folds that keep each group whole. Each dataset is grouped to match what ``unseen'' means in its field. The curation of each dataset, with counts before and after each step, is given in Appendix~\ref{app:datasets}; the representations each dataset uses are listed in Table~\ref{tab:dataset_reps}.

% TODO: review, needs manual review
Different biological domains have established different conventions for what constitutes a good prediction, and applying a single metric everywhere would misrepresent performance in at least some of them. To this end, we assign each task a primary metric matched to its domain: RMSE for molecular property prediction, Spearman rank correlation for antibody regression and developability, AUROC for binary classification (including antibody classification and tumor subtype), weighted F1 for multiclass classification, concordance index for drug-target affinity, and $R^2$ for reaction yield.
% TODO: REVIEW MODIFICATION
Different biological domains have established different conventions for what constitutes a good prediction, and applying a single metric everywhere would misrepresent performance in at least some of them. Each dataset's metric is given in Table~\ref{tab:datasets}, and the reason for each choice is given in ``Models and metrics''.

% TODO: review, needs manual review
No random splits are used. Random splitting inflates performance estimates for molecular data by placing structurally similar compounds on both sides of the train/test boundary \citep{Yang2019, Wallach2018}, and position-level leakage in deep mutational scanning data allows models to memorize site-specific effects rather than learning generalisable sequence-fitness relationships \citep{Notin2023, Vieira2026}. Splits therefore match domain conventions and use 5-fold GroupKFold throughout.
% TODO: REVIEW MODIFICATION
No split is random (Chapter~\ref{chap:background}, ``Evaluation methodology''): regression datasets use five-fold group $k$-fold, and classification datasets use five-fold stratified group $k$-fold.

% TODO: review, needs manual review
\begin{table}[htbp]
\centering
\small
\setlength{\tabcolsep}{3pt}
\caption[The 16 datasets used in the hybrid experiments.]{The 16 datasets used in the hybrid experiments. One row is one dataset after curation. $n$ is the number of samples; the last column gives what the folds are grouped by and the number of groups. $\downarrow$: lower is better; $\uparrow$: higher is better.}
\label{tab:datasets}
\begin{adjustbox}{max width=\textwidth}
\begin{tabular}{P{3.0cm}P{3.3cm}P{2.3cm}P{1.1cm}P{3.5cm}}
\toprule
Dataset & Label (unit) & Task, metric & $n$ & Folds grouped by (groups) \\
\midrule
\multicolumn{5}{l}{\textit{Small molecules: ADME and toxicity}} \\
hERG $K_i$ \citep{Zdrazil2023} & p$K_i$ & regression, RMSE $\downarrow$ & 1,412 & scaffold (769) \\
LogD \citep{openadmet} & LogD (log units) & regression, RMSE $\downarrow$ & 5,028 & scaffold (2,413) \\
Caco-2 efflux \citep{openadmet} & log$_{10}$ efflux ratio & regression, RMSE $\downarrow$ & 2,158 & scaffold (1,086) \\
Rat liver microsome clearance \citep{fang2023, polaris} & log$_{10}$ CL$_\text{int}$ (mL/min/kg) & regression, RMSE $\downarrow$ & 3,053 & scaffold (2,625) \\
MDR1--MDCK efflux \citep{fang2023, polaris} & log$_{10}$ efflux ratio & regression, RMSE $\downarrow$ & 2,641 & scaffold (2,318) \\
\midrule
\multicolumn{5}{l}{\textit{Quantum chemistry}} \\
QM9 \citep{Ramakrishnan2014, qm9_dataset} & HOMO--LUMO gap (eV) & regression, RMSE $\downarrow$ & 9,978 & scaffold (3,452) \\
\midrule
\multicolumn{5}{l}{\textit{Chemical reactions}} \\
Buchwald--Hartwig yield \citep{Ahneman2018} & yield (\%) & regression, RMSE $\downarrow$ & 3,955 & additive (22) \\
Buchwald--Hartwig plates \citep{Ahneman2018} & yield (\%) & regression, RMSE $\downarrow$ & 4,312 & additive (24) \\
\midrule
\multicolumn{5}{l}{\textit{Antibodies}} \\
Koenig expression \citep{Koenig2017, Chungyoun2024} & expression enrichment ratio (no unit) & regression, Spearman $\uparrow$ & 4,275 & mutated position modulo 5 (5) \\
Nanobody thermostability \citep{ValdesTresanco2023, Chungyoun2024} & melting temperature ($^\circ$C) & regression, Spearman $\uparrow$ & 672 & 95\% identity cluster (309) \\
Humanness \citep{Prihoda2022} & human or not (198 : 366) & binary, AUROC $\uparrow$ & 564 & 95\% identity cluster (461) \\
Trastuzumab--HER2 binding \citep{Chinery2024} & binder or not (1,639 : 3,361) & binary, AUROC $\uparrow$ & 5,000 & mutations from wild type (10) \\
\midrule
\multicolumn{5}{l}{\textit{Nucleotide sequences}} \\
H3K4me3 \citep{DallaTorre2025} & mark present or not (9,122 : 9,122) & binary, MCC $\uparrow$ & 18,244 & chromosome (23) \\
TATA promoter \citep{DallaTorre2025} & promoter or not (2,637 : 2,637) & binary, MCC $\uparrow$ & 5,274 & chromosome (23) \\
\midrule
\multicolumn{5}{l}{\textit{Cancer genomics (TCGA)}} \\
Lung adenocarcinoma vs squamous \citep{TCGA2014LUAD, TCGA2012LUSC, Goldman2020} & tumor type (515 : 502) & binary, AUROC $\uparrow$ & 1,017 & tissue source site (68) \\
Breast cancer PAM50 subtype \citep{TCGA2012BRCA, Parker2009, Goldman2020} & subtype (421 : 192 : 141 : 67) & 4 classes, weighted F1 $\uparrow$ & 821 & tissue source site (19) \\
\bottomrule
\end{tabular}
\end{adjustbox}
\end{table}
% TODO: REVIEW MODIFICATION
\begin{table}[htbp]
\centering
\footnotesize
\setlength{\tabcolsep}{3pt}
\caption[The 16 datasets and the best single representation on each.]{The 16 datasets used in the hybrid experiments. One row is one dataset after curation. $n$ is the number of samples; the fifth column gives what the folds are grouped by and the number of groups. The last column gives the best single base representation and its mean score over the five outer folds, in the dataset's own metric. $\downarrow$: lower is better; $\uparrow$: higher is better.}
\label{tab:datasets_modified} % label renamed so the two copies do not clash; keep one table
\begin{adjustbox}{max width=\textwidth}
\begin{tabular}{P{2.6cm}P{2.5cm}P{2.0cm}P{1.0cm}P{2.4cm}P{2.4cm}}
\toprule
Dataset & Label (unit) & Task, metric & $n$ & Folds grouped by (groups) & Best single representation (score) \\
\midrule
\multicolumn{6}{l}{\textit{Small molecules: ADME and toxicity}} \\
hERG $K_i$ \citep{Zdrazil2023} & p$K_i$ & regression, RMSE $\downarrow$ & 1,412 & scaffold (769) & \textbf{[RESULT NEEDED: name, score]} \\
LogD \citep{openadmet} & LogD (log units) & regression, RMSE $\downarrow$ & 5,028 & scaffold (2,413) & \textbf{[RESULT NEEDED: name, score]} \\
Caco-2 efflux \citep{openadmet} & log$_{10}$ efflux ratio & regression, RMSE $\downarrow$ & 2,158 & scaffold (1,086) & \textbf{[RESULT NEEDED: name, score]} \\
Rat liver microsome clearance \citep{fang2023, polaris} & log$_{10}$ CL$_\text{int}$ (mL/min/kg) & regression, RMSE $\downarrow$ & 3,053 & scaffold (2,625) & \textbf{[RESULT NEEDED: name, score]} \\
MDR1--MDCK efflux \citep{fang2023, polaris} & log$_{10}$ efflux ratio & regression, RMSE $\downarrow$ & 2,641 & scaffold (2,318) & \textbf{[RESULT NEEDED: name, score]} \\
\midrule
\multicolumn{6}{l}{\textit{Quantum chemistry}} \\
QM9 \citep{Ramakrishnan2014, qm9_dataset} & HOMO--LUMO gap (eV) & regression, RMSE $\downarrow$ & 9,978 & scaffold (3,452) & \textbf{[RESULT NEEDED: name, score]} \\
\midrule
\multicolumn{6}{l}{\textit{Chemical reactions}} \\
Buchwald--Hartwig yield \citep{Ahneman2018} & yield (\%) & regression, RMSE $\downarrow$ & 3,955 & additive (22) & \textbf{[RESULT NEEDED: name, score]} \\
Buchwald--Hartwig plates \citep{Ahneman2018} & yield (\%) & regression, RMSE $\downarrow$ & 4,312 & additive (24) & \textbf{[RESULT NEEDED: name, score]} \\
\midrule
\multicolumn{6}{l}{\textit{Antibodies}} \\
Koenig expression \citep{Koenig2017, Chungyoun2024} & expression enrichment ratio (no unit) & regression, Spearman $\uparrow$ & 4,275 & mutated position modulo 5 (5) & \textbf{[RESULT NEEDED: name, score]} \\
Nanobody thermostability \citep{ValdesTresanco2023, Chungyoun2024} & melting temperature ($^\circ$C) & regression, Spearman $\uparrow$ & 672 & 95\% identity cluster (309) & \textbf{[RESULT NEEDED: name, score]} \\
Humanness \citep{Prihoda2022} & human or not (198 : 366) & binary, AUROC $\uparrow$ & 564 & 95\% identity cluster (461) & \textbf{[RESULT NEEDED: name, score]} \\
Trastuzumab--HER2 binding \citep{Chinery2024} & binder or not (1,639 : 3,361) & binary, AUROC $\uparrow$ & 5,000 & mutations from wild type (10) & \textbf{[RESULT NEEDED: name, score]} \\
\midrule
\multicolumn{6}{l}{\textit{Nucleotide sequences}} \\
H3K4me3 \citep{DallaTorre2025} & mark present or not (9,122 : 9,122) & binary, MCC $\uparrow$ & 18,244 & chromosome (23) & \textbf{[RESULT NEEDED: name, score]} \\
TATA promoter \citep{DallaTorre2025} & promoter or not (2,637 : 2,637) & binary, MCC $\uparrow$ & 5,274 & chromosome (23) & \textbf{[RESULT NEEDED: name, score]} \\
\midrule
\multicolumn{6}{l}{\textit{Cancer genomics (TCGA)}} \\
Lung adenocarcinoma vs squamous \citep{TCGA2014LUAD, TCGA2012LUSC, Goldman2020} & tumor type (515 : 502) & binary, AUROC $\uparrow$ & 1,017 & tissue source site (68) & \textbf{[RESULT NEEDED: name, score]} \\
Breast cancer PAM50 subtype \citep{TCGA2012BRCA, Parker2009, Goldman2020} & subtype (421 : 192 : 141 : 67) & 4 classes, weighted F1 $\uparrow$ & 821 & tissue source site (19) & \textbf{[RESULT NEEDED: name, score]} \\
\bottomrule
\end{tabular}
\end{adjustbox}
\end{table}

% TODO: review, needs manual review
\paragraph{Pre-processing.} The steps shared across datasets are listed below. The per-dataset steps and the counts after each one are in Appendix~\ref{app:datasets}.
\begin{enumerate}[leftmargin=*]
    \item Every molecule is parsed with RDKit \citep{rdkit}, and molecules that fail are dropped. On hERG $K_i$, LogD and Caco-2 the structures are also standardized: largest fragment kept, charges neutralized, canonical tautomer, and molecules containing a metal removed.
    \item Repeat measurements of one molecule are merged into their median (hERG $K_i$, LogD, Caco-2). hERG $K_i$ molecules whose repeat measurements have a standard deviation above 1.0 p$K_i$ units are removed. $K_i$ values combined from different sources are a known source of large label noise \citep{landrum2024}.
    \item Duplicate molecules are removed by InChIKey on the six small-molecule datasets, keeping the first.
    \item LogD and Caco-2 labels are clipped at the 1st and 99th percentiles of the dataset.
    \item Two datasets are subsampled once with a fixed seed: 10,000 of the 129,428 parsable QM9 molecules, and 5,000 of the 524,346 trastuzumab variants, drawn in proportion to each combination of label and number of mutations.
\end{enumerate}

% TODO: review, needs manual review
\paragraph{Splitting rules.} Four rules hold for every dataset.
\begin{enumerate}[leftmargin=*]
    \item Folds are formed as in Chapter~\ref{chap:background}, ``Evaluation methodology'': no group is split, regression datasets use group $k$-fold, and classification datasets use stratified group $k$-fold.
    \item There are five outer folds, and the stratified folds use a fixed seed.
    \item No split is ever random. What happens when a dataset has fewer groups than folds is stated in Chapter~\ref{chap:background}, ``Evaluation methodology''. Every dataset here has at least five groups, so all outer splits have five folds.
    \item On a given dataset, every representation, model and hybrid is scored on the same five folds.
\end{enumerate}

% TODO: review, needs manual review
The hybrid construction step uses a nested internal train/validation split within each training fold to score representations for feature allocation. This internal split is group-aware, taking the first fold of a grouped 5-fold split (GroupKFold) with the same grouping variable as the outer cross-validation, to prevent the allocation evaluator from seeing related samples on both sides of its split.
% TODO: REVIEW MODIFICATION
Inside each outer training fold, the search scores representations on one group-aware internal split, described in ``From base representations to a hybrid''.

% TODO: review, needs manual review
\paragraph{Groups by domain.} The small-molecule and QM9 datasets are grouped by Bemis--Murcko scaffold, defined in Chapter~\ref{chap:background}, ``Evaluation methodology''. In this chapter the scaffold keeps stereochemistry, and all acyclic molecules share one empty scaffold. On QM9 this acyclic group holds 1,041 of the 9,978 molecules and is the largest group. On the two Biogen datasets (rat liver microsome clearance and MDR1--MDCK efflux) there are 1.2 and 1.1 molecules per scaffold (the $n$ and group columns of Table~\ref{tab:datasets}). On the Biogen datasets, holding out a scaffold is therefore close to holding out single molecules. Chapter~\ref{chap:robustness} used a different scaffold key and different curation for its versions of LogD, Caco-2, hERG $K_i$ and QM9; Appendix~\ref{app:datasets} lists the differences.

% TODO: review, needs manual review
The reaction datasets are Buchwald--Hartwig C--N couplings run on high-throughput plates, in which aryl halide, ligand, base and additive are varied. They are grouped by additive, following the out-of-sample tests of \citet{Ahneman2018}, so each held-out fold contains four or five additives the model never saw in training. The two sets are not independent: they share 3,955 reactions (Appendix~\ref{app:datasets}).

% TODO: review, needs manual review
On the Koenig deep mutational scan of the anti-VEGF antibody G6, every variant carries one mutation. Variants are grouped by the position of that mutation modulo 5, as in ProteinGym \citep{Notin2023}, so each fold is one residue class and holds 836 to 874 variants. All mutations at one position then fall in the same fold, and a model cannot learn a per-position effect and look it up at test time \citep{Vieira2026}. The nanobody and humanness sets are grouped by single-linkage clusters at 95\% sequence identity, after \citet{Kenlay2024}. The identity measure, the cluster sizes and the humanness labels are given in Appendix~\ref{app:datasets}. The trastuzumab variants differ from the wild type only in the ten residues of CDR-H3, and are grouped by how many of those residues are mutated. A held-out fold therefore contains variants at mutation counts not seen in training.

% TODO: review, needs manual review
The two nucleotide tasks come from the revised Nucleotide Transformer benchmark \citep{DallaTorre2025}. Both are grouped by chromosome, so a held-out fold contains whole chromosomes the model never saw in training: between two and six of the 23 chromosome groups (Appendix~\ref{app:datasets} gives the task definitions and how the released files are regrouped).

% TODO: review, needs manual review
The TCGA datasets use RNA-seq expression and other measurements of primary tumors, one sample per patient, obtained through UCSC Xena \citep{Goldman2020}. Samples are grouped by tissue source site, the institution that supplied the sample, which is read from the TCGA barcode. A model therefore cannot score well by recognizing the site a sample came from, and no patient spans training and test.

% TODO: review, needs manual review
\paragraph{Consequences of the grouping.} Group definitions differ between domains on purpose; the one shared rule is that no group spans training and test. Some groupings leave few or very unequal groups, and the folds inherit that. The trastuzumab held-out folds range from 13 to 3,035 variants, with binder shares from 15.5\% to 92.3\% against 32.8\% overall. The breast cancer folds range from 3 to 403 tumors. None of the 68 lung sites supplied both tumor types, so the lung folds are unbalanced: one holds 20 adenocarcinomas and 113 squamous tumors. The Koenig folds are fixed by its five groups. Per-fold scores on these datasets spread more, so results are read per dataset, in the dataset's own metric, and never pooled across datasets.
\subsection{Representations by domain}
% [notes moved to thesis_notes_moved.md]
% TODO: review, needs manual review

We combine a variety of component, or base representations to create these hybrids.

Any representation that gives one fixed-length vector per sample can serve as a base representation, as set out in the description of how the hybrid is built earlier in this chapter. The first version of the method used three molecular representations: PDV, ECFP4 and MHG-GNN. These three are now one small-molecule example among many. Every base representation is frozen: it is computed once from the raw input (a SMILES string, a sequence, a reaction or a patient profile) and is never trained on the labels.

For every data set in this study, we pre-selected ten base representations that spanned across different domains of information and showed some positive results when used on their own for the given prediction task.
% TODO: REVIEW MODIFICATION
For every dataset in this study, we pre-selected at most ten base representations. They span different domains of information, and each showed some positive results on its own for the given prediction task.

To choose them, every candidate representation was first scored on its own with a gradient-boosted tree model \citep{ke2017lightgbm} under the dataset's grouped folds. Candidates were then compared by the Spearman correlation between their out-of-fold predictions, since two representations whose predictions correlate highly carry much the same signal. At most ten were then chosen by hand, looking at each candidate's own score next to these correlations and favoring representations that encode the input in different ways. Two nucleotide datasets have fewer. H3K4me3 had 12 candidates, and four of them predicted almost the same as a better-scoring candidate, with a Spearman correlation of 0.95 or above between their out-of-fold predictions; eight remain. TATA promoter had 10 candidates, one was dropped by the same rule, and nine remain. Augmentation features do not count toward the ten. They are added after the hybrid is built (see the Augmentation subsection above), and the same representation can be a base representation on one dataset and augmentation on another.
% CODE INPUT
On each dataset, the two base representations whose out-of-fold predictions correlate most are \textbf{[RESULT NEEDED: most alike pair, with its Spearman correlation]}. The two that correlate least are \textbf{[RESULT NEEDED: least alike pair, with its Spearman correlation]}.

If an input cannot be parsed or embedded, it gets an all-zero row, so that every representation of a dataset keeps the same rows. In computed descriptors, a missing or infinite value is set to 0, and values are stored unscaled. A pretrained model gives one vector per input, the mean of its last-layer token vectors with padding excluded, unless the tables say otherwise. Scaling and filling of missing values happen inside each fold, on the training rows only, as described in the construction of the hybrid.

\paragraph{Small molecules.} The small-molecule representations (Table~\ref{tab:reps_smallmol}) are the fingerprint, descriptor, graph and language-model families described in Chapter~\ref{chap:background}. The only additions are four sets of 3D shape and pharmacophore features, used as augmentation only. All four are computed from a single conformer per molecule, embedded with a fixed seed and optimized with the MMFF force field. A single conformer may misdescribe flexible molecules.

\paragraph{Chemical reactions.} Reaction datasets (Table~\ref{tab:reps_reactions}) use three kinds of representation. The condensed graph of reaction merges reactants and products into one graph in which the bonds that change are marked \citep{Hoonakker2011}; here it is hashed into a circular fingerprint after atom mapping \citep{Schwaller2021rxnmapper}. Reaction language models read the reaction SMILES directly \citep{Schwaller2021rxnfp, Sagawa2023ReactionT5}. Any molecular representation can also be applied to the reactants, to the product, or to their difference. A one-hot encoding of which aryl halide, ligand, base and additive was used carries no chemistry. It is kept because on plate data such identity features can predict yield about as well as chemical descriptors \citep{Chuang2018}.

\paragraph{Antibodies.} Antibody representations (Table~\ref{tab:reps_antibodies}) start from the heavy- and light-chain amino-acid sequences. Besides one-hot and BLOSUM62 encodings \citep{Henikoff1992}, most are frozen language models, either trained on all proteins or on antibody repertoires only. Most models embed each chain separately and join the two vectors, so their width doubles on datasets with both chains. Several representations read one model in different ways, for instance averaging over earlier layers or over the CDR residues only, with CDRs found by IMGT numbering \citep{Lefranc2003} using ANARCI \citep{Dunbar2016ANARCI}. Datasets of point variants of one parent antibody also use features that describe each mutation.

\paragraph{Nucleotide sequences.} A $k$-mer frequency vector counts how often each DNA word of length $k$ occurs in a sequence, divided by the total number of words \citep{Vinga2003}; the canonical variant merges each word with its reverse complement \citep{Teeling2004}. DNA language models are networks pretrained on genomes to predict masked or next tokens \citep{Zhou2023DNABERT2, DallaTorre2025, Nguyen2023HyenaDNA}. Table~\ref{tab:reps_nucleotides} lists both kinds, along with the small sequence summaries used as augmentation.

\paragraph{Cancer patients.} For the two TCGA datasets (Table~\ref{tab:reps_tcga}), each representation is a different measurement of the same patient's tumor, not a different encoding of one input. The measurements are RNA-seq gene expression (read through foundation models or summarized per chromosome arm), DNA methylation on the HM450 array, gene-level copy number, microRNA abundance, protein levels from reverse-phase protein arrays, and embeddings of the diagnostic whole-slide image. Data were obtained from TCGA through UCSC Xena \citep{TCGA2013, Goldman2020}. The most-variable selections, the z-scoring and the principal components for methylation, copy number, microRNA and protein were fitted once on all patients in a cohort, held-out patients included. No labels are used in these steps. A patient with no measurement for one data type gets an all-zero row for that representation.

We list all of these representations alongside their associated data sets and prediction tasks below:

\begin{longtable}{P{3.0cm}P{5.0cm}P{1.5cm}P{1.5cm}P{1.4cm}}
\caption[Representations used for small molecules.]{Representations used for small molecules. One row is one representation. Role: B = base representation, A = augmentation, B/A = base on some datasets and augmentation on others.}\label{tab:reps_smallmol}\\
\toprule
Representation & What it encodes & Features & Source & Role \\
\midrule
\endfirsthead
\toprule
Representation & What it encodes & Features & Source & Role \\
\midrule
\endhead
\bottomrule
\endfoot
ECFP4 & hashed circular substructures, radius 2 & 2048 & \citep{rogers2010} & B \\
Atom pairs & hashed atom-type pairs with bond distance & 2048 & \citep{Carhart1985} & B \\
Topological torsions & hashed paths of four bonded atoms & 2048 & \citep{Nilakantan1987} & B/A \\
Avalon & hashed paths, rings and atom environments & 2048 & \citep{Gedeck2006} & B \\
MACCS keys & 166 predefined substructure keys & 167 & \citep{Durant2002} & B/A \\
WL graph kernel & Weisfeiler--Lehman label counts, 5 rounds, labels present in at least 5\% of molecules & 140--317 & \citep{Shervashidze2011} & B \\
PDV & RDKit 2D physicochemical descriptors & 200 & \citep{rdkit} & B \\
Mordred & 2D molecular descriptors & 1613 & \citep{Moriwaki2018} & B \\
GROVER & frozen pretrained graph transformer & 4800 & \citep{Rong2020GROVER} & B \\
MHG-GNN & frozen pretrained graph encoder & 1024 & \citep{kishimoto2023} & B \\
GraphMVP & frozen 2D graph encoder pretrained with 3D views & 300 & \citep{Liu2022GraphMVP} & B \\
MoLFormer & frozen SMILES language model; public checkpoint trained on 10\% of its corpus & 768 & \citep{ross2022} & B \\
SMI-TED & frozen SMILES encoder--decoder & 768 & \citep{Soares2025} & B \\
ChemBERT & frozen SMILES language model; community checkpoint trained on ChEMBL following the cited method & 256 & \citep{Zhang2022ChemBERT} & B \\
ChemBERTa-2 & frozen SMILES language model & 384 & \citep{Ahmad2022} & B \\
SELFormer & frozen SELFIES language model & 768 & \citep{Yuksel2023} & B \\
Mol2vec & summed substructure word vectors & 300 & \citep{Jaeger2018} & B \\
3D pharmacophore & counts of 7 pharmacophore types, and mean and spread of distances for the 21 type pairs, one conformer & 49 & \citep{rdkit} & A \\
USR & shape moments from 4 reference points, one conformer & 12 & \citep{Ballester2007} & A \\
USRCAT & USR per pharmacophore atom type, one conformer & 60 & \citep{Schreyer2012} & A \\
Coulomb matrix & sorted eigenvalues of the Coulomb matrix, one conformer & 50 & \citep{Rupp2012} & A \\
\end{longtable}

\begin{longtable}{P{3.0cm}P{5.0cm}P{1.5cm}P{1.5cm}P{1.4cm}}
\caption[Representations used for chemical reactions.]{Representations used for chemical reactions. One row is one representation. Role: B = base representation, A = augmentation, B/A = base on some datasets and augmentation on others.}\label{tab:reps_reactions}\\
\toprule
Representation & What it encodes & Features & Source & Role \\
\midrule
\endfirsthead
\toprule
Representation & What it encodes & Features & Source & Role \\
\midrule
\endhead
\bottomrule
\endfoot
CGR fingerprint & circular fingerprint, radius 2, hashed from the condensed graph of reaction & 2048 & \citep{Hoonakker2011} & B \\
rxnfp & frozen reaction language model & 256 & \citep{Schwaller2021rxnfp} & B \\
ReactionT5 & frozen reaction language model & 768 & \citep{Sagawa2023ReactionT5} & B \\
Component one-hot & which aryl halide, ligand, base and additive & 45 & \citep{Ahneman2018} & B \\
GROVER (reactants) & GROVER averaged over reactants & 4800 & \citep{Rong2020GROVER} & B \\
Mordred (reactants) & Mordred averaged over reactants & 1613 & \citep{Moriwaki2018} & B \\
Mol2vec (reactants) & Mol2vec averaged over reactants & 300 & \citep{Jaeger2018} & B \\
ECFP4 (reactants) & ECFP4, bitwise OR over reactants & 2048 & \citep{rogers2010} & B \\
Avalon (difference) & product minus reactant Avalon bits & 2048 & \citep{Gedeck2006} & B \\
MHG-GNN (difference) & product minus reactant embedding & 1024 & \citep{kishimoto2023} & B \\
MACCS (difference) & product minus reactant MACCS keys & 167 & \citep{Durant2002} & B \\
Mol2vec (difference) & product minus reactant Mol2vec & 300 & \citep{Jaeger2018} & B \\
MACCS (reactants) & MACCS keys, bitwise OR over reactants & 167 & \citep{Durant2002} & A \\
MACCS (product) & MACCS keys of the product & 167 & \citep{Durant2002} & A \\
3D pharmacophore (reactants) & 3D pharmacophore averaged over reactants & 49 & \citep{rdkit} & A \\
3D pharmacophore (product) & 3D pharmacophore of the product & 49 & \citep{rdkit} & A \\
\end{longtable}

\begin{longtable}{P{2.8cm}P{4.6cm}P{2.0cm}P{1.5cm}P{1.4cm}}
\caption[Representations used for antibodies.]{Representations used for antibodies. One row is one representation. LM = language model. Widths marked /chain double on datasets with heavy and light chains (Koenig expression, Trastuzumab--HER2) and are single on the nanobody and heavy-chain-only humanness datasets. -- in Source: a standard encoding or one built for this thesis. Role: B = base representation, A = augmentation, B/A = base on some datasets and augmentation on others.}\label{tab:reps_antibodies}\\
\toprule
Representation & What it encodes & Features & Source & Role \\
\midrule
\endfirsthead
\toprule
Representation & What it encodes & Features & Source & Role \\
\midrule
\endhead
\bottomrule
\endfoot
One-hot & one-hot residues, both chains padded to fixed length & 5600 & -- & B \\
BLOSUM62 & BLOSUM62 row per residue, both chains padded & 5600 & \citep{Henikoff1992} & B \\
Amino-acid composition & fraction of each amino acid, per chain & 40 & -- & A \\
ESM-2 & protein LM, last layer, mean over residues & 1280/chain & \citep{Lin2023} & B \\
ESM-2 (joined chains) & ESM-2 on heavy and light chain joined into one string & 1280 & \citep{Lin2023} & B \\
ESM-2 early layers & ESM-2, mean over layers 0--16 & 1280/chain & \citep{Lin2023} & B \\
ESM-2 middle layers & ESM-2, mean over layers 20--33 & 1280/chain & \citep{Lin2023} & B \\
ESM-2 CDR & ESM-2, mean over CDR residues only & 1280/chain & \citep{Lin2023} & B \\
ESM-2 per CDR & ESM-2, CDR1/2/3 pooled separately then joined & 3840/chain & \citep{Lin2023} & B \\
ESM-2 variable positions & ESM-2, mean over the 20 highest-entropy positions & 1280/chain & \citep{Lin2023} & B \\
ESM-2 shift from parent & ESM-2 variant minus parent embedding & 1280/chain & \citep{Lin2023} & B \\
ESM-2 log-likelihood ratio & masked mutant vs parent log-probability: sum, mean, min & 3 & \citep{Meier2021} & A \\
ESM3 & protein LM & 1536/chain & \citep{Hayes2025} & B \\
ESM3 CDR & ESM3, mean over CDR residues & 1536/chain & \citep{Hayes2025} & B \\
ESM C & protein LM & 1152/chain & \citep{Candido2026} & B \\
ProtT5 & protein LM & 1024/chain & \citep{Elnaggar2022} & B \\
ProtT5 CDR & ProtT5, mean over CDR residues & 1024/chain & \citep{Elnaggar2022} & B \\
SaProt & structure-aware protein LM & 1280/chain & \citep{Su2023SaProt} & B \\
AbLang-2 & paired antibody LM, both chains read together & 480 & \citep{Olsen2024} & B \\
IgT5 & paired antibody LM, both chains read together & 1024 & \citep{Kenlay2024IgT5} & B \\
AntiBERTy & antibody LM & 512/chain & \citep{Ruffolo2021AntiBERTy} & B \\
AntiBERTy CDR & AntiBERTy, mean over CDR residues & 512/chain & \citep{Ruffolo2021AntiBERTy} & B \\
Ab-RoBERTa & antibody LM & 768/chain & \citep{Huh2025} & B \\
Ab-RoBERTa CDR & Ab-RoBERTa, mean over CDR residues & 768/chain & \citep{Huh2025} & B \\
Mutation features & wild-type and mutant residue, substitution score, changes in hydrophobicity, weight and charge, position; averaged over mutations & 48 & -- & B/A \\
\end{longtable}

\begin{longtable}{P{3.0cm}P{5.0cm}P{1.5cm}P{1.5cm}P{1.4cm}}
\caption[Representations used for nucleotide sequences.]{Representations used for nucleotide sequences. One row is one representation. LM = language model. -- in Source: a simple summary built for this thesis. Role: B = base representation, A = augmentation, B/A = base on some datasets and augmentation on others.}\label{tab:reps_nucleotides}\\
\toprule
Representation & What it encodes & Features & Source & Role \\
\midrule
\endfirsthead
\toprule
Representation & What it encodes & Features & Source & Role \\
\midrule
\endhead
\bottomrule
\endfoot
4-mer counts & normalized counts of every 4-letter word & 256 & \citep{Vinga2003} & B \\
5-mer counts & normalized counts of every 5-letter word & 1024 & \citep{Vinga2003} & B \\
6-mer counts & normalized counts of every 6-letter word & 4096 & \citep{Vinga2003} & B \\
Canonical TNF & 4-mer counts merged with reverse complements & 136 & \citep{Teeling2004} & B \\
DNABERT-2 & DNA LM & 768 & \citep{Zhou2023DNABERT2} & B \\
Nucleotide Transformer v2 250M & DNA LM & 768 & \citep{DallaTorre2025} & B \\
Nucleotide Transformer v2 500M & DNA LM & 1024 & \citep{DallaTorre2025} & B \\
HyenaDNA & long-context DNA LM & 256 & \citep{Nguyen2023HyenaDNA} & B \\
GROVER (DNA) & DNA LM (not the molecular GROVER) & 768 & \citep{Sanabria2024} & B \\
Mistral-DNA & mixture-of-experts DNA LM & 768 & \citep{Mourad2024} & B \\
GC content & GC fraction, GC and AT skew, mean and spread of GC over 100-base windows & 5 & -- & A \\
Dinucleotide frequency & frequency of each 2-letter word & 16 & \citep{Vinga2003} & A \\
Codon usage & relative usage of each of the 61 sense codons, read in frame from the first base & 61 & \citep{Sharp1986} & A \\
Sequence complexity & entropy of 1-, 2- and 3-letter words, linguistic complexity, fraction of alphabet used & 5 & -- & A \\
DNA shape & minor groove width, propeller twist, helix twist and roll from a dinucleotide lookup table; mean and spread & 8 & -- & A \\
\end{longtable}

\begin{longtable}{P{3.0cm}P{5.0cm}P{1.5cm}P{1.5cm}P{1.4cm}}
\caption[Representations used for cancer patients (TCGA).]{Representations used for cancer patients (TCGA). One row is one representation. -- in Source: computed for this thesis from the patient's own data. Role: B = base representation, A = augmentation, B/A = base on some datasets and augmentation on others.}\label{tab:reps_tcga}\\
\toprule
Representation & What it encodes & Features & Source & Role \\
\midrule
\endfirsthead
\toprule
Representation & What it encodes & Features & Source & Role \\
\midrule
\endhead
\bottomrule
\endfoot
Geneformer (2,048 genes) & single-cell foundation model (V2-316M) on the 2,048 top-ranked genes, [CLS] vector & 1152 & \citep{Theodoris2023, Geneformer2024card} & B \\
Geneformer (4,096 genes) & the same model on the 4,096 top-ranked genes, [CLS] vector & 1152 & \citep{Theodoris2023, Geneformer2024card} & B \\
BulkRNABert & bulk RNA-seq foundation model & 256 & \citep{Gelard2024} & B \\
Top-variance methylation probes & beta values of the 5,000 most variable CpG probes & 5000 & \citep{Goldman2020} & B \\
Methylation PCA & 50 principal components of the standardized CpG probes above median variance & 50 & \citep{Goldman2020} & B \\
Copy number PCA & 50 principal components of gene-level copy number & 50 & \citep{Mermel2011} & B \\
miRNA & 200 most variable miRNAs, z-scored & 200 & \citep{Goldman2020} & B \\
Protein (RPPA) & reverse-phase protein array, z-scored & 258 & \citep{Li2013TCPA} & B \\
Chromosome-arm expression & mean expression of the genes on each chromosome arm & 42 & -- & B/A \\
TITAN & whole-slide image foundation model & 768 & \citep{Ding2024TITAN} & B \\
Prov-GigaPath & whole-slide image foundation model & 512 & \citep{Xu2024GigaPath} & B \\
CTransPath & histopathology image encoder & 512 & \citep{Wang2022CTransPath} & B \\
Phikon & histopathology image encoder & 512 & \citep{Filiot2023} & B \\
H-optimus-0 & histopathology image encoder & 512 & \citep{Saillard2024} & B \\
HRD score & homologous recombination deficiency: LOH, TAI, LST and their sum & 4 & \citep{Knijnenburg2018} & A \\
Tumor mutation burden & non-silent, missense and truncating mutation counts, log count, rate per megabase & 5 & \citep{Ellrott2018} & A \\
Tumor purity & purity, ploidy, genome doublings, subclonal genome fraction & 4 & \citep{Carter2012} & A \\
Expression statistics & summaries of the patient's expression distribution & 15 & -- & A \\
Immune signatures & 10 CIBERSORT cell fractions and Th1 and Th2 signature scores & 12 & \citep{Thorsson2018} & A \\
CIBERSORT & immune cell fractions deconvolved from bulk expression & 21 & \citep{Newman2015, Thorsson2018} & A \\
Stemness index & stemness indices from expression and from epigenetically regulated expression & 2 & \citep{Malta2018} & A \\
Hypermethylation summary & fraction of all methylation probes with beta above 0.3, and mean beta; a genome-wide stand-in for CIMP status \citep{Toyota1999} & 2 & -- & A \\
\end{longtable}

\begin{longtable}{P{2.7cm}P{2.7cm}P{4.6cm}P{3.1cm}}
\caption[Datasets and the representations each one uses.]{Datasets and the representations each one uses. One row is one dataset. Base representations are the (up to) ten the hybrid search selects features from; augmentation features are added afterward. Datasets marked \textsuperscript{\textdagger} are searched at the fast tier (Table~\ref{tab:search_tiers}). Sample counts, metrics and folds are given in the Datasets and splitting subsection.}\label{tab:dataset_reps}\\
\toprule
Dataset & Task & Base representations & Augmentation \\
\midrule
\endfirsthead
\toprule
Dataset & Task & Base representations & Augmentation \\
\midrule
\endhead
\bottomrule
\endfoot
\multicolumn{4}{l}{\textit{Small molecules}} \\
Caco-2 efflux \citep{openadmet} & efflux ratio (log), regression & WL graph kernel, MHG-GNN, ECFP4, PDV, GROVER, Avalon, Topological torsions, SMI-TED, Atom pairs, MoLFormer & MACCS keys, 3D pharmacophore, USRCAT, Coulomb matrix, USR \\
LogD\textsuperscript{\textdagger} \citep{openadmet} & LogD, regression & MHG-GNN, Mordred, GROVER, ECFP4, SMI-TED, MoLFormer, Avalon, Topological torsions, ChemBERT, GraphMVP & MACCS keys, 3D pharmacophore, USRCAT \\
hERG Ki \citep{Zdrazil2023} & hERG inhibition, pKi, regression & MHG-GNN, Mordred, WL graph kernel, Avalon, GROVER, PDV, ECFP4, Atom pairs, ChemBERT, SMI-TED & MACCS keys, 3D pharmacophore, Coulomb matrix, USR, USRCAT \\
MDR1-MDCK efflux \citep{fang2023} & efflux ratio (log), regression & Mordred, PDV, MHG-GNN, Avalon, GROVER, SMI-TED, Atom pairs, WL graph kernel, MACCS keys, ECFP4 & 3D pharmacophore, Coulomb matrix \\
Rat liver microsome clearance \citep{fang2023} & intrinsic clearance (log), regression & PDV, MHG-GNN, SMI-TED, Mordred, MoLFormer, GROVER, WL graph kernel, Atom pairs, Avalon, ECFP4 & 3D pharmacophore, MACCS keys, USRCAT, Coulomb matrix, USR \\
QM9\textsuperscript{\textdagger} \citep{Ramakrishnan2014} & HOMO--LUMO gap, regression & Mordred, Avalon, SMI-TED, ECFP4, SELFormer, ChemBERTa-2, Mol2vec, ChemBERT, GraphMVP, MoLFormer & Topological torsions, 3D pharmacophore, USRCAT, Coulomb matrix \\
\midrule
\multicolumn{4}{l}{\textit{Chemical reactions}} \\
Buchwald--Hartwig yield \citep{Ahneman2018} & yield (\%), regression & CGR fingerprint, GROVER (reactants), Avalon (difference), Mordred (reactants), MHG-GNN (difference), ReactionT5, Component one-hot, Mol2vec (reactants), rxnfp, MACCS (difference) & MACCS (product), MACCS (reactants), 3D pharmacophore (product), 3D pharmacophore (reactants) \\
Buchwald--Hartwig plates \citep{Ahneman2018} & yield (\%), regression & ECFP4 (reactants), GROVER (reactants), MHG-GNN (difference), CGR fingerprint, Component one-hot, Avalon (difference), ReactionT5, Mol2vec (difference), Mordred (reactants), rxnfp & MACCS (product), MACCS (reactants), 3D pharmacophore (product), 3D pharmacophore (reactants) \\
\midrule
\multicolumn{4}{l}{\textit{Antibodies}} \\
Koenig expression\textsuperscript{\textdagger} \citep{Koenig2017} & expression enrichment of variants, regression & ProtT5, ESM3 CDR, SaProt, IgT5, ESM-2 early layers, ESM C, AbLang-2, ESM-2 (joined chains), ESM-2, ESM-2 middle layers & Mutation features, ESM-2 log-likelihood ratio, Amino-acid composition \\
Nanobody thermostability \citep{ValdesTresanco2023} & melting temperature, regression & SaProt, ESM3, ESM C, ESM-2 variable positions, ESM-2 CDR, ESM-2 early layers, ESM-2 shift from parent, ProtT5 CDR, IgT5, One-hot & Mutation features \\
Humanness \citep{Prihoda2022} & human vs non-human, binary & IgT5, AntiBERTy CDR, Ab-RoBERTa, AbLang-2, BLOSUM62, ESM C, ESM-2 per CDR, ESM-2 early layers, ProtT5, ESM-2 middle layers & none \\
Trastuzumab--HER2 binding\textsuperscript{\textdagger} \citep{Chinery2024} & binder vs non-binder, binary & BLOSUM62, ESM-2 per CDR, ESM-2 middle layers, ProtT5, ESM3, ESM C, Ab-RoBERTa CDR, AntiBERTy, Mutation features, AbLang-2 & none \\
\midrule
\multicolumn{4}{l}{\textit{Nucleotide sequences}} \\
H3K4me3\textsuperscript{\textdagger} \citep{DallaTorre2025} & histone mark present vs absent, binary & GROVER (DNA), Mistral-DNA, Canonical TNF, Nucleotide Transformer v2 250M, Nucleotide Transformer v2 500M, HyenaDNA, 6-mer counts, DNABERT-2 & Codon usage, Sequence complexity, Dinucleotide frequency, DNA shape, GC content \\
TATA promoter\textsuperscript{\textdagger} \citep{DallaTorre2025} & promoter vs non-promoter, binary & Mistral-DNA, GROVER (DNA), HyenaDNA, 6-mer counts, 5-mer counts, Canonical TNF, 4-mer counts, Nucleotide Transformer v2 500M, Nucleotide Transformer v2 250M & Codon usage, Sequence complexity, Dinucleotide frequency, DNA shape, GC content \\
\midrule
\multicolumn{4}{l}{\textit{Cancer patients (TCGA)}} \\
BRCA PAM50 subtype \citep{TCGA2012BRCA} & breast cancer subtype, 4 classes & Geneformer (4,096 genes), BulkRNABert, Protein (RPPA), Copy number PCA, TITAN, H-optimus-0, Prov-GigaPath, CTransPath, Phikon, miRNA & HRD score, Tumor mutation burden, Tumor purity, Expression statistics, Chromosome-arm expression, Immune signatures \\
LUAD vs LUSC \citep{TCGA2014LUAD} & lung cancer type, binary & Geneformer (2,048 genes), Top-variance methylation probes, BulkRNABert, TITAN, Methylation PCA, Chromosome-arm expression, Copy number PCA, miRNA, CTransPath, Protein (RPPA) & Immune signatures, Expression statistics, Stemness index, CIBERSORT, Hypermethylation summary, Tumor mutation burden, HRD score, Tumor purity \\
\end{longtable}

\subsection{Models and metrics}
% [notes moved to thesis_notes_moved.md]

% TODO: review, needs manual review
For a fair comparison, the downstream model pool must cover a range of inductive biases, from methods that degrade in high dimensions (kNN, SVM, GP) to tree ensembles that tolerate irrelevant features. Eleven models are available: tree ensembles (random forest, XGBoost, LightGBM, CatBoost, NGBoost), kernel methods (SVM, Gaussian process), nearest-neighbor (kNN), and neural networks (three-layer DNN, MC-Dropout BNN for uncertainty quantification, and a multi-scale 1D CNN). Standard experiments use five of these: random forest, XGBoost, LightGBM, SVM and DNN, with a Gaussian process added on datasets under 1,000 molecules. We use the full set for model-matching and noise experiments, and skip models with cubic or higher scaling for large datasets: Gaussian processes above 1,000 samples, SVMs above 10,000. No per-dataset tuning is performed: all models use fixed hyperparameters throughout, so performance differences reflect the input, not the model configuration.
% TODO: REVIEW MODIFICATION
For a fair comparison, the downstream models must cover a range of inductive biases. They run from methods that degrade in high dimensions (SVM, GP) to tree ensembles that tolerate irrelevant features. Standard experiments use five downstream models: random forest, XGBoost, LightGBM, SVM and DNN, with a Gaussian process added on humanness. An experiment that uses a model outside these five names it in its own methods. We skip the two models whose fitting cost grows fastest with the number of molecules on large datasets: Gaussian processes above 1,000 samples, SVMs above 10,000. No per-dataset tuning is performed: all models use fixed hyperparameters throughout, so performance differences reflect the input, not the model configuration.

% CODE INPUT
The five downstream models cover bagged trees (random forest), boosted trees (XGBoost and LightGBM), a kernel method (SVM) and a neural network (DNN). We use five because the best set of features depends on the model that uses it. A hybrid that helped only one model family would be a property of that model, not a better input. Tree ensembles are the baseline to beat. A hybrid is a tabular feature matrix, and tree-based models still lead on medium-sized tabular data of about 10,000 samples \citep{Grinsztajn2022}. A Gaussian process \citep{Rasmussen2006} is added as a sixth model on datasets under 1,000 molecules. Three of the 16 datasets are that small: humanness (564 sequences), nanobody thermostability (672 sequences) and breast cancer PAM50 subtype (821 tumors). The Gaussian process runs on humanness only. Humanness is a classification task, so the Gaussian process there has an RBF kernel. The regression version, with an RBF kernel plus a white-noise term, runs on none of the 16 datasets. PAM50 subtype is left out under every augmentation setting for run time. In one search, the one with augmentation features added on top of the budget, the Gaussian process took 44 of its 60 hours. Nanobody thermostability is left out under augmentation on top of the budget, where its searches did not finish in the time limit. A dataset that drops the Gaussian process under any setting gets it under none, so nanobody thermostability is left out everywhere.

% CODE INPUT
Every model uses the same settings on every dataset, so differences between scores come from the input. There is no per-dataset tuning by default. An optional per-fold random search exists. It tries 20 draws of a model's settings, scored in the dataset's own metric on 3 group-aware inner folds of the outer training fold only. The DNN is never tuned this way; it relies on early stopping. \textbf{[RESULT NEEDED: whether tuning pays on base representations, tuned against untuned models, in each dataset's metric]}. Two models are skipped above a cap on the number of molecules. The Gaussian process is skipped above 1,000 molecules, because it builds and factorizes the full molecule-by-molecule kernel. The SVM is skipped above 10,000 molecules. Above that, its kernel cache no longer holds the support vectors, and a fit stops finishing in useful time. Under augmentation on top of the budget, the SVM is skipped above 5,000 molecules. The caps count the molecules the search is given. The search run once per dataset is given the whole dataset. A search inside one outer fold is given that fold's training molecules, about 80\% of the dataset. On two datasets the two counts give different downstream models, both under augmentation on top of the budget. LogD has 5,028 molecules and TATA promoter 5,274. The SVM is dropped when the whole dataset is counted, and kept when about 4,000 training molecules are counted. Some default models are also not run on named datasets where a single fit was too slow. They are listed in Additional file~\ref{app:add:models_left_out}. Every downstream model's settings, library, molecule cap and citation are in Table~\ref{tab:add:model_settings} (Additional files).

% CODE INPUT
$k$-nearest neighbors is not a default downstream model. Where it is used, its default is $k = 5$ neighbors, lowered to the square root of the number of training molecules when a fold has fewer than 25. It is the model in the comparison of importance-weighted against unweighted distances, which reads 5, 10 and 20 neighbors. We state its settings once, here, because Chapter~\ref{chap:discussion}, ``Simple is often better'', relies on that result.

% CODE INPUT
Each dataset is scored in the metric its field normally uses, so its numbers can be compared with published work on the same dataset. Small-molecule, quantum-chemistry and reaction-yield regression use RMSE, which is in the label's own units. Antibody regression uses Spearman $\rho$. These assays are used to rank candidates, and their scales differ between labs. Binary classification uses AUROC, which does not depend on a decision threshold. The two nucleotide sets are the exception: they use MCC. MCC is high only when all four entries of the confusion matrix are good \citep{Chicco2020}. The four PAM50 subtypes, whose class sizes run from 67 to 421 tumors, use support-weighted F1. Table~\ref{tab:datasets} gives the metric of each dataset. The definitions are in Chapter~\ref{chap:background}, ``Evaluation methodology''. Metrics are never averaged across datasets or domains.

% CODE INPUT
Every score is computed on an outer test fold the search never saw. We report the per-fold scores and their spread. Each hybrid is paired by fold with the best single representation on the same fold. Each comparison is reported as its per-fold differences and their spread across the 5 folds, with no threshold and no test. We use no significance tests and no confidence intervals, within or across datasets. With 5 folds, a $p$-value or an interval would rest on 5 correlated values. The fold scores are correlated because the folds share training molecules. No unbiased estimator of the variance of $k$-fold cross-validation exists, and estimates that ignore the overlap underestimate it \citep{Bengio2004}.

% CODE INPUT
A full search takes \textbf{[RESULT NEEDED: hours per dataset for the full search]}. The searches ran on \textbf{[RESULT NEEDED: hardware, as CPU cores, memory, and a GPU where the DNN or the embeddings used one]}.


\subsection{Comparators, controls and ablations}
% [notes moved to thesis_notes_moved.md]
% - TODO(lit): read the Turkina2025 methods for whether high-level fusion pools the predictions of one model per fingerprint, as the ensemble here does.

% CODE INPUT
The hybrid picks the best of many configurations: an allocation method, an importance method, an internal evaluator and a budget. A reader will ask whether its gain is only the reward for searching hard. The experiments in this subsection answer that, one question each. The first five ask whether the search overfits, whether a larger search helps, and whether fusing features beats fusing predictions. They also compare the hybrid with standard feature reduction, and ask whether the internal evaluator must match the final model. The next three measure how stable the selected features are across folds, how the gain changes with dataset size, and whether importance helps during training. The last ones add a few features to one representation, test robustness to label noise, and join two whole representations. All use the group folds of ``Datasets and splitting''.

% CODE INPUT
The hybrid's score is the best of many configurations, while each single representation gives one score. The best of many cross-validated scores is biased upward \citep{Varma2006}. Several experiments therefore use leave-one-fold-out reselection. For held-out fold $f$, we pick the hybrid configuration with the best mean score on the other folds,
\[
c^*_{-f} = \arg\max_c \frac{1}{F-1}\sum_{g\neq f} s_{c,g},
\]
where $s_{c,g}$ is the score of configuration $c$ on fold $g$ and $F$ is the number of folds (arg min for RMSE). We pick the best single representation, $k^*_{-f}$, the same way. We then report
\[
\Delta_f = s_{c^*_{-f},f} - s_{k^*_{-f},f},
\]
signed so that $\Delta_f>0$ means the hybrid is better; for RMSE the sign is flipped. This is arithmetic on stored per-fold scores, so it needs no extra model fits. Importance scores and the feature selection are always recomputed on each fold's training molecules only. Some experiments choose the configuration by a search over all molecules, and redo only the feature selection per fold. Their scores are not fully held out, and we say so where it applies.

% CODE INPUT
The first experiment asks whether the search overfits as the budget grows. It uses nested cross-validation. For every budget $B$, that budget's winning configuration is rebuilt on the outer training molecules and scored on the outer held-out fold. The selection bias at a budget is the score the inner search reported minus the held-out score. It is signed so that a positive value means the search promised more than the held-out fold gave. The budgets are the standard list, 10 to 2,000 features, plus 4,000 and 8,000 features. We add budgets at 0.5, 0.75, 1, 1.25, 1.5 and 2 times the smallest outer-fold training-set size, in molecules. These place a budget where the number of features equals the number of molecules. For the five downstream models, the budget does not set the number of fitted parameters. The DNN and the random forest have more fitted parameters than training molecules at every budget. The SVM is fitted with at most one coefficient per training molecule, so it never can. The capacity of the boosted trees is set by the number and depth of their trees. For these five, the curves therefore say whether the search's choice of budget overfits. This experiment adds one model outside the five, a linear model with almost no penalty. It is ridge regression with penalty $10^{-6}$, or logistic regression with $C = 10^{6}$ on classification. It has one weight per selected feature, so its parameter count equals the budget. We report two numbers per model: the budget each outer fold's search picked, and the budget with the best held-out score. A pick far above the best held-out budget means the search overfits. \textbf{[RESULT NEEDED: held-out score and selection bias against budget, per model and dataset]}.

% CODE INPUT
The size of a search is the number of configurations it scores: budgets times allocation methods, importance methods, internal evaluators and representation counts. At the full tier that is 480 configurations for each pair of downstream model and internal evaluator, and at the fast tier 48 (Additional file~\ref{app:add:search_size}). The question is whether the hybrid's lead over each comparator grows with the number of configurations. The comparators are the ensemble, the best single representation, PCA, and feature reduction at the same width. They get no search. So part of any gap that grows with search size is the hybrid taking the best of more draws. Leave-one-fold-out reselection charges the hybrid for this, and we report the reselected numbers. In the full runs each dataset is searched at one tier only, full or fast. A search-size comparison therefore needs the same dataset searched at two sizes. \textbf{[RESULT NEEDED: hybrid minus each comparator, per dataset, at each search size]}.

% CODE INPUT
The ensemble fits one model per representation and pools their predictions. The hybrid fits one model on one fused feature vector. Joining features before one model is early fusion; pooling the predictions of one model per feature set is late fusion \citep{Stahlschmidt2022}. The hybrid runs its own search on each outer fold's training molecules. Every ensemble choice is made on an inner 3-fold group split of the same training molecules, and the held-out fold is used once. The ensemble either fits the same model on every representation, or lets each representation use its own best model. Its members are either the top 3 representations, or every representation with at least 50 features. The pooling is one of seven: the mean, the median, a 20\% trimmed mean, and two means after removing outliers. The outliers are found by interquartile range or by median absolute deviation. The last two are a mean weighted by each member's inner score, and the mean of the top 3 members. Classification ensembles pool predicted probabilities. Where the metric reads class labels, a majority vote is also offered. There are three controls. The first is the best single representation. The second joins all representations and reduces them by PCA to the hybrid's width. The third is the best pair of model and representation, read off the held-out fold. It is an upper bound, not a fair comparator. Beside each score we report the cost: representations computed at prediction time, models deployed, and fit and prediction time in seconds. The hybrid only saves cost if it draws on fewer representations than the ensemble. An ensemble is the obvious simpler alternative. If it matches the hybrid at lower cost, that result belongs in Chapter~\ref{chap:discussion}, ``Simple is often better''. In QSAR, \citet{Turkina2025} combined six non-hashed fingerprints by low-, mid- and high-level fusion, on six regression and classification datasets. Mid-level fusion, which keeps fingerprint bits by their importance in the per-fingerprint models, gave the best accuracy. \textbf{[RESULT NEEDED: held-out score and cost per method]}.

% CODE INPUT
The next experiment asks whether the gain comes from drawing on several representations, or only from having fewer features. Every method gets the same number of features, $B$, drawn from all base representations and fitted on the fold's training molecules. It runs on regression datasets only. The methods are the hybrid, PCA, mRMR \citep{Peng2005}, recursive feature elimination \citep{Guyon2002}, and $B$ features drawn at random. The hybrid uses its stored configuration at budget $B$, with the features re-selected on each fold. For PCA, all base representations are joined, and each feature is z-scored on the fold's training molecules. PCA is then fitted on those molecules. The number of components is $B$, capped at the total number of features and at one fewer than the number of training molecules. Held-out molecules are projected with the fitted PCA. The PCA control of the ensemble comparison is built the same way. The mRMR ranking stops at 200 features for run time. When the representations hold more than 2,000 features in total, mRMR first keeps the 2,000 with the highest mutual information with the label. Recursive feature elimination uses a random forest of 50 trees of depth at most 5. One downstream model scores every method, a random forest of 100 trees, so only the features differ. This comparison is not fully held out. The hybrid's configuration is chosen using all molecules, and only its feature selection is redone per fold. \textbf{[RESULT NEEDED: held-out score against $B$ (features) for each method]}.

% CODE INPUT
The search scores each representation with an internal evaluator to split the budget. The winning hybrid then goes to a downstream model. This experiment asks whether the two must match. It uses nested cross-validation as in the overfitting experiment, with one search per outer fold for each evaluator setting. Each setting is compared with the matched evaluator. The matched evaluator is the same model family as the downstream model, at the lighter evaluator settings of Table~\ref{tab:add:model_settings}. The near-family setting scores the random forest with LightGBM and the reverse, and XGBoost with LightGBM. It scores the SVM with a Gaussian process, and the DNN with a narrower DNN (64 and 32 units). The far-family setting scores the three tree models with an SVM, the SVM with LightGBM, and the DNN with $k$-nearest neighbors. A fourth setting scores the DNN with a wider DNN (512, 256 and 128 units). The last uses a fixed set of three evaluators, random forest, LightGBM and XGBoost, and keeps the best. We report each setting minus matched, paired per outer fold, at the same number of features. Evaluators with the exact settings of the downstream random forest and LightGBM are implemented. No search uses them, so the full runs do not report them. The answer sets whether one cheap evaluator is enough, or whether the search must be re-run for each model. \textbf{[RESULT NEEDED: each setting minus matched, per downstream model]}.

% CODE INPUT
Selection on each of the $M$ folds gives a binary vector $\mathbf{z}_m \in \{0,1\}^d$, where $d$ is the number of features in all base representations joined. We measure how stable the selection is with the index of \citet{Nogueira2018},
\[
\hat\Phi = 1 - \frac{\frac1d\sum_{f=1}^d s_f^2}{\frac{\bar k}{d}\left(1-\frac{\bar k}{d}\right)}, \qquad s_f^2 = \frac{M}{M-1}\hat p_f(1-\hat p_f).
\]
Here $\hat p_f$ is the fraction of folds that selected feature $f$, and $\bar k$ is the mean number of features selected. $\hat\Phi=1$ means the same features were chosen on every fold. $\hat\Phi \approx 0$ is what random selection gives. $d$ must be that full width, not the union of the selected features. Unstable feature sets would weaken any chemical reading of the blueprints. \textbf{[RESULT NEEDED: $\hat\Phi$ per dataset and model]}.

% CODE INPUT
The learning curves draw subsamples of 10, 25, 50, 75 and 100\% of the molecules, as whole groups, and never fewer than 50 molecules. Above 5,000 molecules we use absolute sizes of 250, 500, 1,000 and 2,000 molecules instead. Each size is drawn 3 times. The 100\% point is drawn once, since every draw is the same. A full search runs on each subsample, with budgets above the number of training molecules dropped. The gain at each size is the hybrid minus the best single representation, with leave-one-fold-out reselection. Small subsamples search fewer budgets. So part of any rise in gain with size may come from searching more configurations. \textbf{[RESULT NEEDED: gain per dataset, model and size]}.

% CODE INPUT
This experiment asks whether importance helps during training, not only during selection. It uses the hybrid at 100 features with random forest, LightGBM and XGBoost, the three models that accept per-molecule weights. Importance is computed on the training molecules with the importance method the search picked for that hybrid. There are five conditions:
\begin{enumerate}[leftmargin=*]
    \item unweighted;
    \item features multiplied by the square root of their importance, after standardizing, with the weights normalized to average 1;
    \item molecules weighted by how well their label agrees with their 10 nearest neighbors in the rescaled space;
    \item both together;
    \item only the features strictly above the median importance.
\end{enumerate}
The weight of molecule $i$ is $1 - |y_i - \bar y_{N(i)}| / (\max y - \min y)$, clipped to between 0.1 and 1. Here $\bar y_{N(i)}$ is the mean label of its 10 nearest neighbors, by Manhattan distance. This comparison is not fully held out. The configuration comes from a search over all molecules, and only the feature selection is redone per fold. The molecule weighting is a noise-mitigation method: it down-weights suspect labels. Tree ensembles are unchanged when a feature is multiplied by a positive constant, because their splits depend only on the order of values \citep{Hastie2009}. So the second condition differs from the unweighted one in one way only. Features with zero importance get weight zero, become constant, and drop out. Every other feature is only rescaled, which the three tree models ignore. The rescaling does matter for the third condition, because it changes which molecules are nearest neighbors. Checking a label against its nearest neighbors goes back to editing a training set by its nearest neighbors \citep{Wilson1972}. Here the molecule is down-weighted, not removed. We found no QSAR study that down-weights training molecules whose labels disagree with their nearest neighbors. The nearest prior work is outside QSAR. Nearest-neighbor noise filters adapted to regression labels remove such samples rather than down-weight them \citep{Martin2021}. Learned per-sample training weights for noisy labels come from general machine learning \citep{shu2019metaweightnet}. \textbf{[RESULT NEEDED: one result line]}.

% CODE INPUT
This experiment asks whether a few extra features help one strong single representation. Two decoys are planted among the candidate features. The first is a block of pure Gaussian noise, as wide as the real candidate features combined. The second is near-copies of 64 base features, with added noise at 5\% of each feature's standard deviation. We count how many decoy features are taken on each fold; a working rule takes none. The same extra features are also offered to the hybrid built from several representations. This asks whether extra features do more on one representation or on the hybrid. \textbf{[RESULT NEEDED: one result line, with the decoy features taken per fold]}.

% CODE INPUT
This experiment asks whether hybrids are more robust to label noise than their base representations. Robustness is the normalized area under the retention curve of Chapter~\ref{chap:robustness}, \aucnorm (Equation~\ref{eq:auc_norm}). There are two differences from Chapter~\ref{chap:robustness}. First, the retention curve uses each dataset's own metric, not $R^2$. For an error metric the ratio is inverted, so 1 still means no loss. Second, there are seven noise levels: 0, 0.2, 0.3, 0.5, 0.75, 1.0 and 1.5. Each is a fraction of the spread of the fold's clean training labels. Censoring has its own axis, the fraction of labels clipped: 0, 0.10, 0.20, 0.25, 0.30 and 0.40. We compare the hybrid at each budget with each single representation. We also compare it with the best single representation, re-picked in each fold. One downstream model, a random forest, is used under each regression noise type. The hybrid's settings are searched again at every noise type and level, on the noisy training labels. Held-out labels stay clean. The search uses the first fold's scores to drop configurations, so that fold is left out of both sides and four folds are reported. In Chapter~\ref{chap:robustness} the model, not the representation, explained most of the variance in \aucnorm. Here the model is fixed, so only the representation changes.

% CODE INPUT
The last experiment asks whether joining two whole representations scores above the better of the two, or between them. For each dataset we take its top three single representations and form each of the three pairs. We score each representation alone and each concatenation, with the same model and folds. A concatenation that lands between its two parts supports the case for combining and filtering at once. \textbf{[RESULT NEEDED: for each pair, the score of each representation alone and of their concatenation]}.


\subsection{Transfer across datasets, antibody mutational scans and data held at separate sites}
% [notes moved to thesis_notes_moved.md]
% - TODO(lit): check that \citet{Hanser2025} states this weighting in its methods before citing it for it; otherwise cite the collaborator's released code.

% CODE INPUT
This subsection sets up three studies outside the single-dataset setting. They reuse search settings on a second dataset, build hybrids on antibody mutational scans, and select features with data held at separate sites. We describe how each setup was built and scored, with a short background for each; the outcomes are in the results sections. One requirement covers all three. Each needs several numeric representations of the same objects, one supervised target, and a grouped split that keeps near-duplicates on one side.

% CODE INPUT
The first study asks whether search settings found on a large dataset can be reused on a small one, instead of running a new search. The search settings are the importance method, the allocation method, the budget $B$ and the internal evaluator. They are decisions, not a set of features. The source dataset is LogD (lipophilicity), with 5,028 molecules. The target dataset is hERG binding affinity ($\mathrm{p}K_i$), with 1,412 molecules. The candidate representations are those listed for both datasets, kept only where the width (number of features) matches. Both searches then range over the same features. The current lists share seven representations: Avalon, ECFP4, Mordred, and four pretrained embeddings (GROVER, MHG-GNN, SMI-TED and ChemBERT). The width check can still drop some of them. \textbf{[RESULT NEEDED: the shared representation set the full run uses]}. The search is the expensive step. Settings that transfer would make it a one-off cost, paid on the largest dataset. The closest prior work is Meta-QSAR \citep{Olier2018}. It chose the learning method and molecular representation for a new QSAR dataset from results on over 2,700 other QSAR datasets. It beat random forest on a fingerprint by up to 13\% on average. Like the source search settings here, it reuses decisions, not a set of selected features. We have not found a paper that applies feature indices selected on one QSAR dataset unchanged to another, but our search of the literature was not complete.

% CODE INPUT
All four comparisons run inside every held-out fold of the target. The first is the best single representation, re-chosen on that fold's training molecules. The second is a hybrid searched on that fold's training molecules, which is the upper reference. The third applies the source search settings, with the features re-selected on the target's training molecules. The fourth applies the source's own selected feature indices unchanged, so no target label is used to pick features. The last two separate two questions: whether the search decisions transfer, and whether the chosen features transfer. We use five scaffold-group folds, RMSE in $\mathrm{p}K_i$ units, and the five downstream models (random forest, LightGBM, XGBoost, SVM and DNN). \textbf{[RESULT NEEDED: hybrid minus best single representation for each transfer comparison, per downstream model]}.

% CODE INPUT
A deep mutational scan measures a property, here expression or binding, for many single- or few-point mutants of one antibody. It is used to map which positions tolerate change. We use four scans with a single wild type, taken from FLAb \citep{Chungyoun2024}. Two are the expression and binding scans of the anti-VEGF antibody G6 \citep{Koenig2017}; the expression scan is one of the 16 datasets of this chapter. The other two are the expression and binding scans of the anti-fluorescein antibody 4-4-20. The candidate representations are the protein language-model embeddings and cheaper sequence features of ``Representations by domain''. We add two representations of the change from the wild type. The first has 5 features: the change in hydrophobicity, charge, isoelectric point, molecular weight and aromaticity, summed over the mutations. Each property is z-scored over the 20 amino acids first. The second places the same five features at the mutated position, so it has five features for every sequence position. We use nested cross-validation. The search sees the outer training variants only, and the hybrid and every single representation are scored on the untouched outer fold. The continuous scans are scored by Spearman correlation. The 4-4-20 binding scan is binarized, binder if its $\mathrm{p}K_d$ is above 5.0, and scored by AUROC. We report hybrid minus best single representation, per scan, on the outer folds. The two wild-type-difference representations are also scored as single representations, and each run reports their rank among all representations. The number of features each representation got is stored per fold with every result, but this comparison does not report it. So hybrid minus best single says nothing about the wild-type-difference representations unless that count is also reported. \textbf{[RESULT NEEDED: features the hybrid took from each of the two wild-type-difference representations, per scan, per fold]}. \textbf{[RESULT NEEDED: hybrid minus best single representation per scan, per downstream model]}.

% CODE INPUT
Federated learning trains one model across data held by several parties without moving the data \citep{McMahan2017}. It has been run across pharmaceutical companies on shared property-prediction targets \citep{Heyndrickx2024}. In this field the motivation is commercial. Structures and assay results are the asset, so pooling them is not possible \citep{Hanser2023,Bassani2023}. Our collaborator, Lhasa Limited, uses the teacher-student process described in Chapter~\ref{chap:background}, ``Noise in molecular data'' \citep{Hanser2025}. We use the collaborator's code for the fingerprints, the site models and the site clustering, so the numbers stay comparable to theirs. The step that combines the site models' predictions is re-implemented from that code.

% CODE INPUT
The site models' predictions are combined in one of two ways. Under distance weighting, each site model's vote on a molecule is weighted by how close that molecule is to the site's own training molecules. A site counts most on chemistry it has seen. Closeness is the mean distance between fingerprints, from the query molecule to its 8 nearest neighbors among a random 20\% of the site's training molecules. The distance is the one set for the run, Manhattan by default; the collaborator's code uses Euclidean. For each site, these distances are rescaled to 0--1 over the query molecules. Each is turned into a similarity as $(1 - \text{distance})^2$, and the similarities are normalized to sum to 1 over the sites. Under score weighting, each site instead gets one weight, from its own cross-validated MCC. Only predictions on shared molecules cross a site boundary. In the published process, a student model is then trained on the most confident combined labels, the same number from each class. In the hERG comparison below no student is trained. The combined predictions score the held-out molecules directly.

% CODE INPUT
The addition in this thesis is that feature selection, not only model fitting, happens at the sites. Each site runs the hybrid search on its own molecules. Features are chosen on 80\% of its molecules and scored on the other 20\%, split by scaffold group. Each site reports three things: feature indices, importance values and one score. No molecule and no feature value leaves the site. A coordinator splits the total budget $B$ across sites in proportion to their scores. Each share $b_s$ is rounded down, with at least one feature per site, so $\sum_s b_s \approx B$. If every site scores 0, the budget is split equally. Only the proportional split is run. With no equal-split comparison, we cannot say what the site scores buy. A feature another site has already taken is skipped, and the site takes its next-ranked feature instead, so the features in the hybrid are all distinct. Indices and importances are smaller, and harder to invert, than model weights or gradients. Splitting the budget between sites is also the same question the allocation method answers between representations. In related work, each client runs a mutual-information feature selection on its own data, and a server aggregates the results, with no raw data exchanged \citep{Cassara2022}.

% CODE INPUT
hERG channel blockade is a cardiotoxicity liability screened early in drug discovery, and its public data comes from many labs \citep{Czodrowski2013}. Three real sources act as sites. The first is the public ChEMBL hERG $K_i$ set used elsewhere in this thesis. The other two are the collaborator's training set (7,772 molecules) and test set (4,383 molecules). The two collaborator sets were labeled by different curation rules, so the sites differ in curation, not only in chemistry. This is source heterogeneity that exists in the world. The process is usually simulated with sites made by clustering one dataset. Molecules held by more than one source are deduplicated on InChIKey, keeping the public copy. The public set's target is binary: a molecule is a blocker if $\mathrm{p}K_i \ge 6$, that is $K_i \le 1\,\mu\mathrm{M}$. We use five scaffold-group folds over the pooled molecules. The search is scored by AUROC. The reported number is MCC, which does not reward predicting the majority class.

% CODE INPUT
The comparison crosses two choices: molecules pooled or kept at their sites, and all features or a selected set. Pooled molecules with all features is the upper reference, not possible in practice. Pooled molecules with selected features measure what the selection is worth on its own. Molecules at their sites with all features is federated training as usually done. Molecules at their sites with selected features is the only setting where nothing but indices, importances, one score and predictions moves. Selected minus all features, at fixed site structure, is what selection buys. Pooled minus site-held, at fixed features, is what pooling buys. We add four context comparisons. They are the single best source, picked on training molecules; both ways of weighting the site models; the same search on one fingerprint alone; and one site model per representation. The budget $B$ is swept over 50, 100, 250, 500 and 1,000 features, so no comparison rests on one width. Both all-features settings fit one 2,048-bit Morgan fingerprint of radius 2, with pharmacophoric atom types. Both selected settings fit features drawn from the hERG $K_i$ dataset's candidate representations. So selected minus all features differs in two ways at once, width and source. The budget sweep varies the width of the selected set only. It cannot separate the source difference. The single-fingerprint search is the closest control for the source difference: the same search, on ECFP4 only. Two curators disagreeing on one molecule is label noise with a known source (Chapter~\ref{chap:similarity}).

% CODE INPUT
Four other separate-site experiments are run. In the first, the sites carry different amounts of label noise. Selection at the sites is compared with the site-model ensembles and with one model trained on all sites' molecules pooled, at several noise levels. One version has each site first remove the molecules it flags as noisy. In the second, every site holds the same molecules as a different representation. Fusing features is compared with pooling the site models' predictions. In the third, the sites hold different families of representations (substructure, physicochemical and pretrained) of shared molecules. In the fourth, the sites are antibody sources: a binding-energy-change set, a deep mutational scan of binding, and a thermostability set. Within a site, sub-sites are made by clustering language-model embeddings, not fingerprints. The antibody version skips the two ensemble comparisons. Each site's features are reduced separately, so one site's model cannot predict on another site's features. The labels are also in incompatible units (kcal/mol, a normalized score and degrees Celsius). \textbf{[RESULT NEEDED: one line on the outcome of each of the four]}.

% CODE INPUT
The separate-site work has four limits. One source is proprietary and cannot be redistributed, so the comparison cannot be reproduced from public data alone (Appendix~\ref{app:availability}). The comparison depends on the collaborator's implementation of the process. The published process labels a large separate public set and trains a student on it. Where a student is trained here, the training molecules play that role, and held-out molecules are excluded from it. This is a compromise. Finally, the sites differ in size and in how their labels were derived, so a site's score reflects its curation as well as its chemistry.

% CODE INPUT
Three external benchmark tasks, not among the current datasets, are possible later additions and are not run here. The first is activity-cliff prediction on the dopamine D2 receptor, factor Xa and SARS-CoV-2 main protease, following a published comparison of representations and models \citep{Dablander2023}. The second is antibody developability, on the biophysical properties of a public competition set, using that competition's cluster splits \citep{vanNiekerk2026}. The third is antibody--antigen binding free-energy change, asking how close sequence representations get to a structure-based graph network \citep{Hummer2025}.


\section{Results: when the hybrid helps}
% ---- NOTES: Results: when the hybrid helps ----
% Takeaway: Dataset by dataset, where the hybrid scores better than the best single representation and than an ensemble, and where it does not.
% Takeaway: Where it does not help, one representation may already carry most of the signal; a tie there is a sign the search is not inflating scores.
% Paragraph 1: how to read every number in this section
% - Every comparison is on the outer held-out folds of the group splits.
% - Report the leave-one-fold-out reselected numbers for both the hybrid and the best single representation. Reselection charges both sides for being a maximum.
% - Each dataset is reported in its own metric. Give per-fold scores and their spread. No significance tests, and no averaging across datasets.
% Paragraph 2: hybrid against the best single representation, dataset by dataset
% - One row per dataset and downstream model: the best single representation and its score, the hybrid's score, and hybrid minus best single, fold by fold.
% - Sign every difference so that positive means the hybrid is better, in both metric directions.
% - Group the rows by the six domains: small molecules, the quantum-chemical property, reactions, antibodies, nucleotide sequences, cancer genomics.
% - Report the count of datasets, per model, where the reselected hybrid is ahead on the fold mean. Say what the top and bottom of the count are (datasets ahead, datasets run).
% - The best single representation is chosen per fold, so it can differ between folds. Name it in each fold.
% - Table (\label{tab:hybrid_vs_single}): one row is one dataset scored with one downstream model. Columns: domain; dataset; metric; the best single representation in each fold; its score in each fold; the hybrid's score in each fold; the difference in each fold (positive = hybrid better); mean and spread of the difference across the 5 folds.
% - Result needed: the per-dataset table from the full runs.
% Paragraph 3: hybrid against an ensemble of one model per representation
% - Report three differences per (dataset, model): hybrid minus the ensemble that chooses its own members and pooling rule on the inner split, hybrid minus best single, and the upper bound read off the held-out fold.
% - The upper bound is not a fair comparator. Report it only to show how far the fair comparators are from it.
% - Report the cost beside the score: representations computed at prediction time and models deployed. The hybrid saves only where it draws on fewer representations than the ensemble.
% - If an ensemble matches the hybrid at lower cost on a dataset, say so here; it also belongs in Chapter 7, "Simple is often better".
% - Result needed: hybrid, ensemble and best single per dataset and model, with the cost figures.
% Paragraph 4: how the gain changes with the number of molecules
% - Report hybrid minus best single at each size, reselected, per dataset and model.
% - At the smallest sizes the training folds can hold fewer molecules than the smallest budget. Those sizes search fewer budgets. Say at which sizes this happens, per dataset.
% - So a gain that rises with size can come from more data or from searching more budgets. The curve alone cannot separate the two. State this next to the figure.
% - Figure: x = number of molecules (log scale), y = hybrid minus best single in the dataset's metric, positive = hybrid better; one line per model; one panel per dataset; per-fold points shown.
% - Result needed: gain against number of molecules, per dataset and model.
% Paragraph 5: does a larger search give a larger gain? (only if paired runs at two search sizes are made)
% - A search-size comparison needs the same dataset searched at two numbers of configurations, with folds and representations unchanged. The full runs search each dataset at one tier only.
% - An early comparison of a larger search (5 folds, 16 configurations per model) against a smaller one (3 folds, 2 to 6 configurations) is not usable. It differs in three ways at once, not one: the number of configurations, the number of folds, and the representations searched.
% - If the paired runs are made: report hybrid minus each comparator at the two sizes, reselected, per dataset and model. The two tiers differ in which allocation methods, importance methods and representation counts are searched, not only in how many configurations. Say so next to the numbers.
% Paragraph 6: where the hybrid does not help
% - Name the datasets and models where the reselected hybrid is behind or level with the best single representation.
% - For each, report what share of the hybrid's features comes from the best single representation. That count is stored per fold with every searched configuration, so it needs no extra run.
% - This tests one explanation: where one representation carries most of the signal, the search can only add features from weaker ones. Report the share; do not state the explanation as found.
% - Also report training-fold size and the smallest budget searched, since small datasets search fewer budgets.
% - Result needed: the datasets where the hybrid is not ahead, with the best representation's share of the hybrid's features.
% ---- END NOTES ----
% TODO: Results TODO(prelim): methods are final, numbers will be replaced by the server runs.

% TODO: RESULTS CAVEAT: state the re-run explicitly, once, here, and cross-reference it from ch.6 and ch.7.

\section{Results: why the hybrid helps}
% ---- NOTES: Results: why the hybrid helps ----
% Takeaway: The gain comes from which features are chosen, not only from having fewer: at the same number of features the hybrid is ahead of PCA and random features.
% Takeaway: Different downstream models want different numbers of features and different representations.
% Paragraph 1: what this section asks
% - Five questions, one paragraph each. Does the gain come from drawing on several representations, or only from having fewer features? Must the internal evaluator be the downstream model? Does importance help when used in training? Is adding a few features to one strong representation safe? Do different downstream models want different hybrids?
% - Scores are on held-out group folds, in each dataset's own metric. Give per-fold scores and their spread. No significance tests.
% Paragraph 2: several representations, or only fewer features?
% - If the hybrid leads PCA and random features at the same $B$, the gain is not only from having fewer features.
% - Report the realized number of features for every method at each $B$. Compare only where the widths match. At large $B$ the hybrid can end with fewer than $B$ features.
% - One sentence of caveat: the hybrid's configuration was chosen using all molecules.
% - Preliminary (to be replaced by the full runs), widths matched, PCA, recursive feature elimination and random features only (no mRMR). Antibody thermostability (Tresanco), Spearman, hybrid / PCA / recursive feature elimination / random features:
%   - 10 features: 0.5945 / 0.4860 / 0.4682 / 0.2145.
%   - 50 features: 0.6131 / 0.5225 / 0.5689 / 0.5094.
%   - 250 features: 0.6029 / 0.4824 / 0.6004 / 0.5636.
% - Same runs, hERG Ki at 10 features: hybrid 0.6304 RMSE (pKi) against 0.6776 for PCA; the hybrid scored better on all 5 folds.
% - Same runs: at 1,000 features the methods ended at different widths, so that budget cannot be compared.
% - Figure: x = $B$ (features, log scale), y = held-out score in the dataset's metric, one line per method, one panel per dataset.
% - Result needed: held-out score against $B$ for each method, per dataset.
% Paragraph 3: must the internal evaluator be the downstream model?
% - Report each evaluator setting minus the matched evaluator, paired per outer fold, at the same number of features.
% - For random forest, LightGBM and XGBoost the three-evaluator set contains the matched evaluator. It cannot score worse on the inner folds than matching.
% - The finding to state, whichever way it goes: which evaluator best predicts what the downstream model will do with a feature set.
% - If the matched evaluator is not always best, read it as informative, not as a failure. The default search tries four internal evaluators and keeps whichever scores best, so it already covers both outcomes.
% - Why the same number of features matters: a difference between settings that ended at different widths is not a difference in evaluator.
% - No numbers exist yet for near family, far family or the wider network.
% - Table: one row per downstream model, per dataset. Columns, each minus matched (mean and spread over folds): near family, far family, wider network (DNN row only), three-evaluator set.
% - Result needed: each setting minus matched, per downstream model and dataset.
% Paragraph 4: does importance help when used during training?
% - If the feature-scaling condition stays, read it as "delete the zero-importance features" on these three tree models.
% - Report how many of the hybrid's features get zero importance, per dataset (features out of 100). This says how much the condition can change.
% - The same scaling used as a distance for nearest-neighbor prediction is reported in Chapter 6, not here.
% - Preliminary, 250-feature hybrid: LightGBM gain gave exactly zero importance to 66 to 78 of the 250 features.
% - Same runs, Caco-2: every weighting tried scored at or below unweighted, for random forest and for LightGBM.
% - Result needed: each condition minus unweighted, per model and dataset, and the zero-importance counts.
% Paragraph 5: adding a few features to one strong representation
% - Base: each of the dataset's ten representations in turn, used whole. Report the best single representation per dataset in the main table.
% - Control at the same width: the same number of extra features drawn at random from the other representations. With-extras minus random-extras separates the selection from simply adding features.
% - Report the score with and without extras, per model and dataset.
% - Report how many of the features taken per fold are real, how many are near-copies of base features, and how many are pure noise.
% - If decoys make up a large share of what is taken, say that a score gain alone does not show the extras carry signal.
% - Report whether extras do more on one representation or on the multi-representation hybrid.
% - Preliminary, antibody thermostability (Tresanco, 672 antibodies), 3 folds, Spearman. Base: SaProt alone (1,000 features). Candidate extras: the 48 engineered mutation features. Up to 20 extras taken; no admission rule.
%   - Random forest: 0.5683 without extras, 0.5762 with them.
%   - SVM: 0.5716 without extras, 0.5830 with them.
%   - Features taken per fold, random forest, 3 folds: real 12, 7, 18; near-copy decoys 21, 11, 21; noise decoys 7, 2, 1.
% - Table: one row per (dataset, model). Columns: base representation, score without extras, score with extras, real features taken, near-copy decoys taken, noise decoys taken (per fold).
% - Result needed: score with and without extras, and decoy counts (features per fold), per dataset and model.
% Paragraph 6: do tree ensembles and neural networks want different hybrids?
% - Two parts: the number of features each model does best with, and which representations each model does best with.
% - Number of features: for each downstream model, the budget with the best held-out score.
% - Representations: each whole base representation is fed to each downstream model, 5 folds. Rank the representations by score for each model. These scores come from the full runs.
% - Report the Spearman correlation of those rankings between each pair of models, per dataset (about 10 representations per dataset).
% - A low correlation between a tree model and the DNN means the two models would fund different representations. Name the best representation per model where they differ.
% - Then compare the features chosen by a tree-model hybrid and a DNN hybrid on the same dataset: overlap as Jaccard, $|A\cap B|/|A\cup B|$, over the selected features, broken down by source representation.
% - The search is run separately for each downstream model, so each model gets its own budget and representations.
% - Preliminary, nanobody thermostability, 5 folds: SVM did best at 300 features (0.6384 Spearman); DNN, LightGBM and random forest at 50 features; XGBoost at 100.
% - Figure: heatmap of rank correlation between downstream models, one panel per dataset.
% - Figure: for one dataset, held-out score against budget (features, log scale), one line per downstream model, showing where each peaks.
% - Result needed: best budget per model and dataset; rank correlation per model pair and dataset; Jaccard overlap per dataset.
% ---- END NOTES ----
% TODO: Results TODO(prelim): methods are final, numbers will be replaced by the server runs.


\section{Results: blueprints on specific datasets}
% ---- NOTES: Results: blueprints on specific datasets ----
% Takeaway: On datasets with named features, the best hybrid's selected features are read against the known chemistry of the endpoint.
% Takeaway: Which settings win differs by dataset, and the selection is only partly stable across folds.
% Paragraph 1: what this section does
% - Take the best hybrid on a few datasets and read its selected features against the known chemistry of the endpoint.
% - Only named features can be read chemically. ECFP bits, MACCS keys and the physicochemical descriptor vector (PDV) are named; Mordred, atom-pair and Avalon features are indices only.
% - So choose datasets whose representations include ECFP4 or PDV: hERG Ki, Caco-2 permeability or lipophilicity.
% - Of each dataset's 10 representations, hERG Ki and Caco-2 carry two readable ones (ECFP4 and PDV). Lipophilicity carries one (ECFP4).
% Paragraph 2: how the features are read
% - Two downstream models are used: random forest and XGBoost.
% - Per fold, each of the two is trained on its own hybrid, built from its own best search result. No model's hybrid is read under another model.
% - SHAP values are computed only for the model that scores best on that fold, on the fold's held-out molecules. Each feature's score is its mean absolute SHAP value.
% - Both models are tree ensembles, so SHAP values come from TreeSHAP \citep{Lundberg2020}.
% - Every selected feature is mapped back to its representation. Absolute SHAP values are summed per representation and scaled to sum to 1: each representation's share of the hybrid.
% - The shares are computed inside each fold, then averaged over folds. The list of top features comes from one fold, because each fold selects different features.
% - An ECFP bit is resolved from the first of the dataset's first 50 molecules that sets it, and that molecule's atom environment is drawn. A bit none of those 50 molecules sets stays unresolved.
% - One hashed bit can stand for several environments, so say which one was drawn.
% Paragraph 3: one worked example, hERG Ki
% - The budget split across representations: how many features each representation contributes, and its share of summed importance.
% - The top named features, with their substructures or descriptor names.
% - Read them against what is known about hERG binding: most blockers carry a basic nitrogen, protonated at physiological pH, and hydrophobic or aromatic groups. In the channel pore these meet two aromatic residues, Tyr652 and Phe656.
% - Sources, neither yet in refs.bib: Aronov, "Predictive in silico modeling for hERG channel blockers", Drug Discovery Today 10(2):149--155, 2005, doi 10.1016/S1359-6446(04)03278-7 (hERG blocker pharmacophore models). Mitcheson et al., "A structural basis for drug-induced long QT syndrome", PNAS 97(22):12329--12333, 2000, doi 10.1073/pnas.210244497 (the Tyr652 and Phe656 binding residues).
% - State for each top feature whether it confirms known chemistry, suggests something not expected, or cannot be read (unnamed feature).
% - Selected means predictive, not causal; correlated features can stand in for one another. One sentence.
% - Result needed: per-representation feature counts and importance shares for the hERG Ki hybrid, for the model explained on each fold.
% - Result needed: the top 10 to 20 named features with their SHAP values.
% Paragraph 4: a second dataset, for contrast
% - Same reading as Paragraph 3, shorter: does the selection pick features known to drive permeability or lipophilicity?
% - Result needed: as for hERG Ki.
% Paragraph 5: how stable the selection is across folds
% - Measured with the Nogueira stability index. 1 means the same features in every fold; about 0 is what random selection gives.
% - Preliminary, 3 folds: 0.447 on Caco-2 permeability, 0.275 on hERG Ki, 0.074 on nanobody thermostability (Tresanco).
% - Each of these runs used 5 of the dataset's 10 representations, and only random forest and SVM as downstream models. For thermostability all 5 were learned embeddings. The numbers will change with the full set.
% - Preliminary: the search was also repeated with one more fold held out, and the result scored on that fold. The score was worse on every fold of all three datasets.
% - Caco-2: RMSE 0.2718 to 0.2953, 0.2986 to 0.3289, 0.2740 to 0.3077 (own fold to one fold fewer). hERG Ki: RMSE in pKi 0.5702 to 0.6392, 0.5862 to 0.6636, 0.5605 to 0.6034.
% - Thermostability: Spearman 0.6335 to 0.5743, 0.6703 to 0.6529, 0.5530 to 0.5286.
% - This comparison changes two things at once: the search sees one fold fewer of molecules, and it produces different features. It does not test whether features chosen on one fold carry over to another fold.
% - What a low index means for the chemical reading: correlated features can swap between folds and carry the same information.
% - Per-fold spread only; no confidence intervals and no significance tests.
% - Table: stability index, one row per dataset, one column per downstream model.
% - Result needed: stability index per dataset and downstream model, 5 folds.
% Paragraph 6: does the blueprint change across chemical space?
% - Molecules are split into 5 k-means clusters of chemical space, fitted on the training molecules only. Importance and allocation are computed separately in each cluster.
% - Question: do different regions give each representation a different share of the budget? For example, do substructure fingerprints take more of the budget in regions with many activity cliffs?
% - Compare the per-representation budget shares between clusters with a Spearman rank correlation, averaged over cluster pairs.
% - Control: the same training molecules split at random into parts of the same sizes. Clusters matter only if their shares agree less than the random parts do.
% - Cluster edges: the clustering is repeated with a fraction of the training molecules (0, 0.01, 0.05, 0.1 and 0.2) moved at random to another cluster. Say whether the per-cluster shares change as the edges blur.
% - Also compare the selected features themselves between clusters: two clusters can split the budget the same way and still take different features.
% - The same question on scaffold-defined subsets of a dataset: the best search settings found in each subset.
% - Result needed: per-cluster budget shares by representation, hERG blockers (Lhasa Limited), against the random-partition control.
% Paragraph 7: do the best search settings co-vary across datasets?
% - For each dataset and downstream model: the winning allocation method, importance method, internal evaluator, representation count and budget.
% - Which datasets land on the same settings, and whether data complexity measures computed on each dataset track the winning budget. One Spearman coefficient per complexity measure and downstream model; coefficients only, no p-values.
% - This is the one place the full-run winners per dataset are reported.
% - Figure: one small panel per setting (allocation method, representation count, importance method, internal evaluator). One bar group per dataset; bar = how many (downstream model, budget) combinations each value won.
% - Figure: for one hERG Ki hybrid. Left: stacked bar of the budget split across representations (features per representation). Right: the top named features, with ECFP and MACCS substructures drawn on an example molecule and descriptor names listed.
% - Table: one row per dataset (hERG Ki and the contrast dataset). Columns: downstream model, budget in features, features per representation, top 5 named features.
% - Result needed: table of winning settings, one row per dataset, one column group per downstream model.
% ---- END NOTES ----
% TODO: Results TODO(prelim): methods are final, numbers will be replaced by the server runs.

\section{Results: transfer across datasets and antibody mutational scans}
% ---- NOTES: Results: transfer across datasets and antibody mutational scans ----
% Takeaway: Search settings found on lipophilicity, with features re-selected on hERG, beat hERG's best single representation on 3 of 5 downstream models; copying the features unchanged did so on 2 of 5 (preliminary).
% Takeaway: On an antibody expression scan the hybrid was ahead of the best single representation on all 5 downstream models (preliminary).
% Paragraph 1: what this section reports
% - First: search settings found on lipophilicity reused on hERG $\mathrm{p}K_i$. Second: the hybrid against the best single representation on an antibody deep mutational scan.
% - All numbers below are preliminary and will be replaced by the full runs.
% Paragraph 2: reusing search settings, the numbers
% - Metric: RMSE in $\mathrm{p}K_i$ units on hERG, lower is better. Five scaffold-group folds.
% - Each difference is the best single representation's RMSE minus the comparison's RMSE, mean over the 5 folds. Positive means the comparison had the lower error.
% - The best single representation is chosen per fold. It was not the same representation on every fold and model.
% - Re-selecting features on hERG with the source search settings beat the best single representation on 3 of 5 downstream models. Example: SVM, 0.0938 lower RMSE.
% - Copying the source's selected features unchanged beat it on 2 of 5 models. Example: SVM, 0.0791 lower RMSE.
% - Against the hybrid searched on hERG itself: re-selecting with the source settings came within 0.01 RMSE of it on 4 of 5 models (SVM 0.5678 against 0.5693). The gap is 0.0121 on DNN, where the source settings scored lower.
% - Give the spread with each number. The fold standard deviation of the copied-features difference on DNN is 0.1823 RMSE, larger than any mean difference in the table.
% - State what the two comparisons separate: whether the search decisions carry over, and whether the chosen features carry over. Let the table carry the numbers; no verdict beyond the counts.
% - Table: one row per downstream model (random forest, XGBoost, SVM, DNN, LightGBM). Columns: best single representation RMSE; hybrid searched on hERG RMSE; source settings re-selected RMSE; source features copied RMSE; difference for each of the last two against the best single, with fold standard deviation. Caption: metric, units, 5 folds, positive means lower error than the best single representation.
%   Source settings, features re-selected: SVM +0.0938, DNN +0.0651, LightGBM +0.0300, random forest -0.0037, XGBoost -0.0222.
%   Source features copied unchanged: SVM +0.0791, LightGBM +0.0314, DNN -0.0222, random forest -0.0203, XGBoost -0.0052.
% - Result needed: the same table from the full runs.
% Paragraph 3: the antibody mutational scan, the numbers
% - Dataset: the expression scan \citep{Koenig2017}, 4,275 antibody variants of one wild type, on average 1.0 mutation per variant.
% - Nested cross-validation with 3 outer folds of unequal size: 874, 1,710 and 1,691 test variants (3,401, 2,565 and 2,584 training). Reduced search: 200 search configurations, budgets of 100 and 500 features.
% - Metric: Spearman correlation, higher is better. Difference is hybrid minus best single representation, mean over the 3 outer folds.
% - The hybrid was higher than the best single representation on all 5 downstream models. Example: LightGBM, +0.0431 Spearman.
% - Per outer fold, the hybrid was higher in 14 of 15 model-fold pairs (5 models, 3 outer folds). The exception is XGBoost on one fold (-0.030 Spearman).
% - The best single representation was ProtT5 on every model and every fold. Its held-out Spearman, mean over the 3 outer folds, ranged from 0.6511 (XGBoost) to 0.7154 (SVM).
% - The full run covers all four scans with the full search and the full list of budgets. The binding scans are scored by AUROC, not Spearman.
% - Table: one row per downstream model. Columns: best single representation (name, Spearman); hybrid Spearman; difference, mean and the 3 per-fold values. Caption: 3 outer folds, nested cross-validation, reduced search.
%   DNN +0.0286, LightGBM +0.0431, random forest +0.0366, SVM +0.0287, XGBoost +0.0257 (hybrid minus ProtT5, Spearman).
% - Result needed: the other three scans, and the full search on the expression scan.
% Paragraph 4: which features the antibody hybrids selected
% - Averaged over the 15 selected hybrids (5 models, 3 outer folds), the budget went to ProtT5 (33% of features), IgT5 (29%), AbLang-2 (24%) and ESM-2 (12%).
% - One-hot and BLOSUM sequence encodings together took 0.9% of features.
% - The two representations of the change from the wild type were selected by 6 of the 15 hybrids. Where selected, the two together took at most 6% of the features.
% - So report this as a hybrid-against-best-single result on antibodies, not as a result about the wild-type-difference representations. One line on them at most.
% - Figure (optional): stacked bar of the budget split across representations for each of the 15 antibody hybrids.
% ---- END NOTES ----
% TODO: Results TODO(prelim): methods are final, numbers will be replaced by the server runs.

\section{Results: data held at separate sites}
% ---- NOTES: Results: data held at separate sites ----
% Takeaway: Does feature selection run separately at each site match feature selection on pooled molecules, when only feature indices, importances, one score per site and predictions cross between sites?
% Paragraph 1: how to read this section
% - The main comparison uses three real hERG sources as sites. The other separate-site experiments get one line each.
% - Every number is MCC on the held-out folds (five scaffold-group folds). Give per-fold scores and their spread. No significance tests.
% - The budget is swept at 50, 100, 250, 500 and 1,000 features. Report every budget, not one.
% Paragraph 2: the hERG two-by-two
% - The four settings: molecules pooled or kept at their sites, crossed with all features or a selected set.
% - Selected minus all features, at fixed site structure: what the selection buys.
% - Pooled minus site-held, at fixed features: what pooling buys. The pooled settings need all molecules in one place, so they are references, not options open to the sites.
% - Report the selection difference separately for pooled and for site-held molecules. Say at which budgets it keeps its sign.
% - Table (\label{tab:sites_two_by_two}): one row per budget (features). Columns: pooled/all, pooled/selected, site-held/all, site-held/selected, selected minus all (pooled), selected minus all (site-held). Mean MCC and spread over the five folds.
% - Figure: the two-by-two, with the MCC of each setting at one budget written in its box. Per-budget numbers stay in the table.
% - Result needed: MCC for the four settings at each budget, per fold.
% Paragraph 3: how to read the selection difference
% - Selected minus all features differs in two ways at once: the number of features and the representations they come from.
% - With pooled molecules, pooled/selected minus the ECFP4-only search, at equal width, is what drawing on several representations buys.
% - With site-held molecules, report the selection difference as a mix of width and source.
% - Result needed: how many representations reach all three sources (6 of 10, preliminary).
% - Result needed: pooled/selected minus the ECFP4-only search, at each budget.
% Paragraph 4: the context comparisons
% - Best single source, picked on its own training molecules and scored on the fold. Its gap to the site-held selected setting is what combining sites is worth.
% - Also report the best single source picked on the held-out fold. It needs the answers, so it is a bound, not a comparator.
% - The collaborator's two teacher-weighting schemes: per-molecule fingerprint distance, and one weight per site from its own cross-validated score.
% - One teacher per representation, predictions weighted by each representation's cross-validated MCC.
% - How the coordinator split the budget: each site's share is in proportion to the score it reported. Report each site's score and its share of the features, per budget.
% - A naive pool (one model on all molecules and the fingerprint) is the pooled/all corner. Do not report it twice.
% - Result needed: MCC of each context comparison at each budget; each site's score and share of the features.
% Paragraph 5: other separate-site experiments, one line each
% - Include those that came out null; say what was measured and what the difference was, in the experiment's own metric.
% - Sites with different amounts of label noise: selection at the sites against the teacher ensembles and a naive pool, at several noise levels. One version has each site remove the molecules it flags as noisy first.
% - Representations as sites: every site holds the same molecules as a different representation. Feature-level fusion against pooling the teachers' predictions.
% - Families of representations as sites (substructure, physicochemical, pretrained): one selection over all families against a selection per family, at the same total budget.
% - Antibody sources as sites: a binding-energy-change set, a deep mutational scan of binding, and a thermostability set. Report per-site scores, and the transfer of one site's selected feature positions to another.
% - The antibody sites' labels are in different units (kcal/mol, a normalized score, degrees Celsius). Two teacher-ensemble comparisons are not defined there and were not run. Say so in one sentence.
% - Copying a feature position between sites assumes the feature means the same at every site. State the antibody transfer result only for representations where that holds.
% - Result needed: one number per experiment, in its own metric, at each budget.
% Paragraph 6: what these results can and cannot support
% - Sites differ in curation as well as chemistry, so a gap between sites is not only a chemistry effect.
% ---- END NOTES ----
% TODO: Results TODO(prelim): methods are final, numbers will be replaced by the server runs.


\section{Discussion}
% [notes moved to thesis_notes_moved.md]
% - TODO(lit): neither source above is a study of when combining gives no gain. A source that reports it directly is still needed.

% TODO: review, needs manual review
Ambroise and McLachlan \citep{Ambroise2002} showed the extreme case of feature selection overfitting: an SVM achieved near-zero cross-validated error on random labels when gene selection was performed on the full dataset before cross-validation, but when selection was nested inside each fold the median error was 11 out of 40, close to the 50\% expected by chance. Vabalas et al.\ \citep{Vabalas2019} showed that the bias introduced by pre-selection can exceed the bias from hyperparameter tuning, and Cawley and Talbot \citep{Cawley2010} argued that feature selection is itself a form of model selection whose complexity cost must be accounted for. KIRBy follows the standard remedy: we compute all importance scores within each cross-validation fold, on training data only. But the problem of ``search degrees of freedom'' \citep{Tibshirani2015} remains. Searching over many possible feature subsets adds complexity even when each individual selection is unbiased. The budget acts as a hard capacity constraint: selecting 100 features from a pool of 5,000 limits the space of possible selections far more tightly than $L_1$ regularization over the same pool, because the feasible set grows polynomially in the budget while the continuous parameter space of a regularized model is unbounded \citep{Ng2004}. Important here is a distinction: the features being selected are frozen embeddings, not parameters being fitted. Choosing which pre-existing knowledge to trust is, we think, a less expressive operation than learning new features, and Kumar et al.\ \citep{Kumar2022} showed that frozen embeddings fed through a linear probe can outperform full fine-tuning by up to 10\% in out-of-distribution accuracy precisely because fine-tuning distorts pretrained structure. And stability can be measured directly: the Nogueira index \citep{Nogueira2018} quantifies the consistency of selected feature sets between folds. If selection were overfit, feature sets would differ wildly from fold to fold.

\section{Limitations}
% [notes moved to thesis_notes_moved.md]
% - TODO(lit): source for SchNet outperforming frozen embeddings on 3D quantum-mechanical properties. The sentence has no citation. The SchNet paper (Sch\"utt et al., NeurIPS 2017) and Faber et al. 2017 do not compare SchNet with frozen or pretrained embeddings.

% TODO: review, needs manual review
We want to be clear about KIRBy's limitations. For small datasets (below about 200 samples), the hybrid search cannot reliably outperform the best single source. We suspect this is because the internal validation set within each fold shrinks to roughly 20 samples, and at that size the evaluator cannot distinguish genuine improvements from noise. Tuning cannot fix this. And at the other extreme, frozen embeddings have a ceiling. No signal can be added beyond what the inputs already contain, and on tasks where the true signal requires end-to-end learning from raw inputs (certain 3D quantum-mechanical properties, for example, where SchNet outperforms any frozen embedding by a wide margin), this ceiling binds. Although the methodology naturally accommodates embeddings from sequence-trained models (ESM-2, ProtBERT, etc.), encoding raw sequences directly as one-hot or $k$-mer features disperses the positional and compositional information that makes sequences useful in the first place. Integrating 20 to 50 foundation models into a single pipeline, each with its own dependencies, GPU requirements, and version constraints, is genuinely difficult engineering. Embeddings need only be computed once per dataset and everything is cached. Yet the initial integration burden is real and scales with the number of models in the pool. Finally, on endpoints with high experimental noise, even a well-constructed hybrid may not improve on simpler baselines. Kramer et al.\ \citep{Kramer2012} documented that experimental uncertainty in heterogeneous public $K_i$ data is approximately 0.54 log units. This sets a natural accuracy ceiling of $R^2 \approx 0.81$ regardless of model sophistication; the budget constraint plausibly guards against overshoot, but it cannot push past the noise floor.


\section{Conclusions}
% ---- NOTES: Conclusions ----
% Takeaway: One paragraph per question this chapter answers, then a short cost line, then the bridge to Chapter 6.
% - Quote each question's wording from Section~\ref{sec:intro:questions} at the start of its paragraph. State the answer in one or two sentences.
% Paragraph 1: how do you select the best set of features for a task? (does the hybrid beat its inputs)
% - Answer against two comparators: the best single representation at its full width, and an ensemble of one model per representation.
% - Framing: where the hybrid only matches the best single representation, the search still found that representation without it being named in advance.
% - The premise tested here, argued in Chapter 4: different representations carry different information.
% - Result needed: hybrid against each comparator, per dataset, in the dataset's own metric, with per-fold spread.
% - Result needed: the datasets where the hybrid did not help, and which single representation dominated on each.
% Paragraph 2: selection, not only reduction
% - At the same number of features, the hybrid is compared against PCA, mRMR, recursive feature elimination and random features.
% - If the hybrid is ahead at matched width, the gain comes from which features are chosen, not from having fewer features.
% - The evidence is the comparison at matched numbers of features, not the information-bottleneck analogy.
% - Result needed: hybrid minus each comparator at matched width, per dataset.
% - Result needed: best budget per downstream model (features), and whether tree ensembles and neural networks prefer different widths.
% - Result needed (only if the search-size comparison is kept): hybrid minus each comparator at a small and a large search, after leave-one-fold-out reselection.
% Paragraph 3: one search setting is not enough
% - Every searched setting can pick a different value on different datasets: allocation method, number of representations, importance method, internal evaluator, downstream model, budget.
% - Consequence: a fixed recipe (for example, top two representations, one importance method, one budget) would lose score on some datasets. This is why the search exists.
% - Framing only, not proof: no single optimizer is best when averaged over all problems \citep{Wolpert1997}.
% - Result needed: per setting, the number of datasets on which each value was chosen.
% Paragraph 4: is the search overfitting?
% - The whole search runs inside each outer fold, and importance is computed on training molecules only (nested cross-validation). One sentence.
% - Result needed: selection bias per budget (features), and the stability index of the selected features per dataset.
% Paragraph 5: what blueprints tell a chemist
% - Descriptor features and unhashed substructure keys have chemical names. Features from learned embeddings do not.
% - Result needed: one or two datasets where the named features that were selected match known chemistry of the endpoint.
% Paragraph 6: beyond small molecules, transfer, and data held at separate sites
% - One sentence on what a dataset needs for the hybrid to help: enough samples, and base representations whose features carry importance.
% - Label noise sets a separate ceiling on any gain (the public $K_i$ uncertainty estimate).
% - Result needed: antibodies, reactions, nucleotides and gene expression: hybrid against the best single representation on each.
% - Result needed: transfer across datasets, the antibody mutational scans, and data held at separate sites. One sentence each on the outcome.
% Paragraph 7: are hybrids more robust to label noise than their base representations? (only if this comparison stays in this chapter)
% - Result needed: \aucnorm, hybrid at each budget against each single representation, per dataset and noise type, with per-fold spread.
% Paragraph 8: does importance weighting help? (the parts this chapter answers)
% - Here weighting means using importance scores to scale features or molecules during training, not only to rank and select features.
% - Five training conditions on the hybrid: unweighted, features scaled by importance, molecules weighted by agreement with their neighbors, both, and only the features above the median importance.
% - Caveat: the configuration comes from a search over all molecules; only the feature selection is redone per fold.
% - Weighting in distances and under learning with noisy labels is answered in Chapter 6.
% - Result needed: each weighted condition minus unweighted, per dataset and model, in the dataset's own metric.
% Paragraph 9: cost, stated once
% - One or two sentences: the search is expensive, and it must be run again for new folds, a new representation or a new downstream model.
% - Simpler feature selection may be enough on some tasks.
% Paragraph 10: bridge to Chapter 6
% - Chapter 4 argued that a fixed fingerprint with Tanimoto similarity spends most of its features on information unrelated to the task.
% - This chapter used the hybrid only as input to a model.
% - Chapter 6 uses it for similarity, in two ways, and asks whether either improves noise detection and learning with noisy labels, which Chapter 4 showed failing under standard similarity.
% - Why it matters: if a feature set chosen for the task gives better neighbors, every method that relies on neighbors could use it. Examples: noise detection, applicability domain, nearest-neighbor prediction.
% - Figure (optional): three boxes in a row. Chapter 4 (standard representations and similarity fail) to Chapter 5 (a feature set chosen for the task) to Chapter 6 (that feature set used for similarity and label noise).
% ---- END NOTES ----
